% =====================================================================
%  LHL-PAPER.TEX -- the two c/0004 notes, combined.
%
%  This file is the union of proof.tex and lhl-bound.tex, which both
%  remain in this folder unchanged and are the editable originals.
%  Nothing mathematical was added, removed or restated here.  What the
%  merge did, and nothing else:
%
%   (1) One preamble.  The macro sets of the two notes are unioned;
%       lhl-bound's \extpub was folded into proof.tex's \extp, so the
%       public-seed advantage has one spelling throughout.
%   (2) One front matter, one \cjresolves, one bibliography (all five
%       entries, alphabetical), one running head per part.
%   (3) Two label collisions resolved.  lhl-bound's \label{fact:cs}
%       (Cauchy-Schwarz, two forms) became fact:csforms, proof.tex
%       keeping fact:cs; lhl-bound's \label{sec:setting} is gone, the
%       setting now being stated once, in Part I.
%   (4) lhl-bound's Section 1 was a verbatim restatement of the game
%       "so this note stands alone".  It is replaced by
%       Section~\ref{sec:budgets}, which keeps only what that section
%       ADDED -- the query budgets r_S, s_S and the collision
%       statistics CP, c_k, cbar -- and points back to Part I for the
%       game.  Likewise lhl-bound's display of the conjecture is gone;
%       \eqref{eq:conj} is now Part I's own display of it, which
%       proof.tex labelled eq:pubconj.
%   (5) Cross-file pointers rewritten.  "proof.tex" -> Part I,
%       "main.tex" -> the game of Section~\ref{sec:setting} or the
%       statement of Section~\ref{sec:posed}, "this note" -> this
%       paper / Part I / Part II as the scope of each sentence
%       requires.  main.tex is still cited by name where the point is
%       the statement note itself.
%   (6) One repair, not a merge artefact.  proof.tex starred "The
%       conjecture, as posed" and then \label'd it, so both of its
%       "Section~\ref{sec:posed}" references expanded to nothing (see
%       proof.aux, \newlabel{sec:posed}{{}{1}...}).  The section is
%       numbered here, which fixes the references and leaves their
%       prose alone.  It is Part I's Section 1.
%
%  Both originals' own header warnings still apply verbatim to the
%  text below; read them in proof.tex and lhl-bound.tex before
%  revising anything here.
% =====================================================================
\documentclass{conjura-solution}

% ---- document-specific macros (union of the two notes')
\newcommand{\Fun}{\mathsf{Fun}}
\newcommand{\Adv}{\mathbf{Adv}}
\newcommand{\Gm}{\mathsf{Game}}
\newcommand{\Pred}{\mathrm{Pred}}
\newcommand{\Ext}{\mathrm{Ext}}
\newcommand{\Extp}{\mathrm{Ext}\text{-}\mathrm{pub}}
\newcommand{\pred}{\mathrm{pred}}
\newcommand{\ext}{\mathrm{ext}}
\newcommand{\extp}{\mathrm{ext}\text{-}\mathrm{pub}}
\newcommand{\sd}{\mathit{sd}}
\renewcommand{\sample}{\stackrel{{\scriptscriptstyle\$}}{\leftarrow}}
\newcommand{\SD}{\mathrm{SD}}
\newcommand{\calK}{\mathcal{K}}
\newcommand{\calD}{\mathcal{D}}
\newcommand{\calR}{\mathcal{R}}
\newcommand{\calS}{\mathcal{S}}
\newcommand{\CP}{\mathrm{CP}}
\newcommand{\ret}{\textbf{return}\ }
\newcommand{\emp}{\mathrm{emp}}

% ---- extra theorem-like environments the class doesn't provide
\newtheorem*{theoremstar}{Theorem}
\newtheorem{fact}{Fact}

% ---- part heads, in the house style; \section is styled by the class
\newcommand{\cjpart}[2]{%
  \clearpage
  {\noindent\sffamily\bfseries\color{cjAccent}\footnotesize PART #1}\par
  \vspace{0.35em}%
  {\noindent\LARGE\itshape #2\par}
  \vspace{0.5em}%
  {\color{cjAccent}\hrule height 0.7pt}%
  \vspace{1.4em}%
}

% ---- authorship/review provenance, per part: the two notes were
% written and refereed separately and the merge did not re-referee them
\newcommand{\AIModelsUsed}{Claude Opus 5 (Anthropic)}
\newcommand{\PromptedBy}{Pooya Farshim}
\newcommand{\AIDisclosure}{%
  \begingroup
  \small\centering
  \begin{tabular}{@{}rl@{}}
    \textbf{AI:} & \AIModelsUsed \\
    \textbf{Prompts:} & \PromptedBy \\
    \textbf{Reviewers:} & Part~I: nobody. Part~II: two blind adversarial referees (AI); no human \\
    \textbf{Date:} & Part~I, 14 August 2026; Part~II, 29 August 2026 \\
  \end{tabular}
  \par
  \endgroup
}

\runninghead{LHL FOR RANDOM FUNCTIONS: PUBLIC SEED}

\begin{document}

\cjkicker{CONJURA \textperiodcentered{} SOLUTION}
\cjtitle{LHL for Random Functions}
\cjsubtitle{The Public-Seed Conjecture, Proved; and What the Leftover Hash Lemma Gives}
\cjresolves{The public-seed conjecture of \emph{Randomness Extraction
from Unpredictable Random-Oracle Sources: A Leftover-Hash-Lemma
Conjecture, Public Seed} (see \texttt{main.tex} in this folder) ---
proven true, with $c = 2$. Part~I proves it. Its secret-seed companion, which no longer has a
statement note of its own and is recorded below, is not addressed and
remains open; a natural but false transposition of it to the
public-seed game is also disproven here
(Proposition~\ref{prop:false}). Both statements are reproduced
verbatim below.
\par\smallskip\noindent Part~II resolves nothing further. It asks one question instead: set beside
the bounds already known for extraction from a random oracle, what does the conjectured bound buy?
It buys seed length, and by a wide margin --- $77$ bits in the leading example against $190$,
$1920$ and $113$.}

\vspace{0.6em}
\AIDisclosure

\vspace{1em}


\cjpart{I}{The Conjecture, Proved}

\section{The conjecture, as posed}\label{sec:posed}

The two conjectures below share a setting and differ only in whether the distinguisher is given
the seed. Conjecture~\ref{conj:stmt-pub} is the whole of \texttt{main.tex} in this folder,
reproduced verbatim; it is the statement this paper resolves.
Conjecture~\ref{conj:stmt-secret} is its secret-seed companion, which had a statement note and a
page of its own until both were removed from the site on 28 August 2026, and which is reproduced
here from that note, verbatim, as the only surviving record of it. It is printed rather than merely
cited because Proposition~\ref{prop:false} is a statement about an expression derived from it, and
Remark~\ref{rem:cost} compares against it. They are numbered here secret seed first, because the
argument below reaches the public-seed statement by way of the secret-seed one; the numbering is
this document's own. That shared setting is summarised first and then given in full, in the form
used throughout the proof, in Section~\ref{sec:setting}.

\medskip
\noindent\textbf{Setting.}
Fix nonempty finite sets $\calK$ (seeds), $\calD$ (inputs), $\calR$ (outputs), and write
$K := |\calK|$, $D := |\calD|$, $R := |\calR|$. Let $\Fun(\calK \times \calD, \calR)$ be the set of
all functions $\calK \times \calD \to \calR$, let $\SD(\cdot,\cdot)$ denote statistical distance and
$U_{\calR}$ the uniform distribution on $\calR$. The source $S$, the predictor $P$ and the
distinguisher $\mathsf{D}$ are computationally unbounded and receive the \emph{entire} function
table of $H$ as an explicit input. In the prediction game $\Gm\ \Pred^{S}_{P}$ the predictor is
to guess the source's output $x$ from $H$ and the source's auxiliary output $z$; in the two
extraction games a seed $\sd \sample \calK$ is drawn after $S$ has terminated and $\mathsf{D}$ is
to tell $H(\sd,x)$ from uniform, given $H$ and $z$. The two extraction games differ only in
whether the seed $\sd$ is kept secret from $\mathsf{D}$ (the game $\Gm\ \Ext$) or given to it
(the game $\Gm\ \Extp$); $\Gm\ \Pred$ and $\Gm\ \Extp$ are written out in
Figure~\ref{fig:games} and Definition~\ref{def:games} below.

\medskip
\noindent\textbf{Advantages.}
Set $\Adv^{\pred}_{\calK,\calD,\calR,S}(P) := \Pr[\Pred^{S}_{P} \Rightarrow 1]$ and
\[
  \Adv^{\ext}_{\calK,\calD,\calR}(S,\mathsf{D}) := 2\Pr[\Ext^{S}_{\mathsf{D}} \Rightarrow 1] - 1 ,
  \qquad
  \Adv^{\extp}_{\calK,\calD,\calR}(S,\mathsf{D}) := 2\Pr[\Extp^{S}_{\mathsf{D}} \Rightarrow 1] - 1 .
\]
A source $S$ is \emph{$\epsilon$-unpredictable} if
$\Adv^{\pred}_{\calK,\calD,\calR,S}(P) \le \epsilon$ for every unbounded predictor $P$.

\begin{conjecture}[Secret seed]\label{conj:stmt-secret}
There is a universal constant $c > 0$, independent of $\calK$, $\calD$, $\calR$ and $\epsilon$, such
that for all nonempty finite $\calK, \calD, \calR$, all $\epsilon \in (0,1]$, every
$\epsilon$-unpredictable source $S$ and every unbounded distinguisher $\mathsf{D}$,
\[
  \Adv^{\ext}_{\calK,\calD,\calR}(S,\mathsf{D})
  \;\le\; \delta(\epsilon, K, D, R)
  \;:=\; c \cdot \sqrt{\frac{\epsilon R + \log_2 D}{K}} \; .
\]
\end{conjecture}

\begin{conjecture}[Public seed]\label{conj:stmt-pub}
There is a universal constant $c > 0$, independent of $\calK$, $\calD$, $\calR$ and $\epsilon$, such
that for all nonempty finite $\calK, \calD, \calR$, all $\epsilon \in (0,1]$, every
$\epsilon$-unpredictable source $S$ and every unbounded distinguisher $\mathsf{D}$,
\[
  \Adv^{\extp}_{\calK,\calD,\calR}(S,\mathsf{D})
  \;\le\; \delta_{\mathrm{pub}}(\epsilon, K, D, R)
  \;:=\; c \cdot \Bigl( \sqrt{\epsilon R} + \sqrt{\frac{\log_2 D}{K}} \,\Bigr) .
\]
\end{conjecture}

\noindent
Conjecture~\ref{conj:stmt-pub} is what this paper settles: it holds with $c = 2$
(Corollary~\ref{cor:pub}), as a consequence of the sharper Theorem~\ref{thm:main}.
Conjecture~\ref{conj:stmt-secret} is a statement about the hidden-seed game $\Gm\ \Ext$, is not
addressed here, and remains open.

\section{Introduction}\label{sec:intro}

Part~I revisits the secret-seed conjecture of Section~\ref{sec:posed} under the one change that
the seed $\sd$ is given to the distinguisher. That game is written $\Gm\ \Extp$ and its advantage
$\Adv^{\extp}$, with the suffix carried everywhere, so that no statement here can be mistaken for
a statement about the hidden-seed game $\Gm\ \Ext$ of the companion notes; both are set out in
Figure~\ref{fig:games} and Definition~\ref{def:games}. The conjecture below is
Conjecture~\ref{conj:stmt-secret}, the secret-seed conjecture, stated there for
$\Adv^{\ext}$ and transposed verbatim to $\Adv^{\extp}$. It is not
Conjecture~\ref{conj:stmt-pub}, the separate and weaker conjecture posed for the
public-seed game; that one is treated in Corollary~\ref{cor:pub} and is not refuted.

\begin{conjecture}[Conjecture~\ref{conj:stmt-secret} transposed to the public-seed game]\label{conj:main}
There is a universal constant $c > 0$, independent of $\calK$, $\calD$, $\calR$ and $\epsilon$, such
that for all nonempty finite $\calK, \calD, \calR$, all $\epsilon \in (0,1]$, every
$\epsilon$-unpredictable source $S$ and every unbounded distinguisher $\mathsf{D}$,
\[
  \Adv^{\extp}_{\calK,\calD,\calR}(S,\mathsf{D})
  \;\le\; \delta(\epsilon,K,D,R)
  \;:=\; c \cdot \sqrt{\frac{\epsilon R + \log_2 D}{K}} \; .
\]
\end{conjecture}

\noindent
Three things happen. Conjecture~\ref{conj:main} becomes \emph{false}, and demonstrably so: a
maximum-entropy source with no auxiliary output already refutes it
(Proposition~\ref{prop:false}). The proof technique of the companion proof note, on the other
hand, survives essentially unchanged, and delivers
\[
  \Adv^{\extp}_{\calK,\calD,\calR}(S,\mathsf{D})
  \;\le\; \frac{1}{\sqrt{2}}\sqrt{\epsilon R}
  \;+\; \min\Bigl\{\, \frac{6}{5}\sqrt{\frac{1 + \ln D}{K}} \;,\;\;
        \sqrt{\frac{\ln(2eD\epsilon)}{2K}} + \frac{9}{10\sqrt{K}} \,\Bigr\}
\]
(Theorem~\ref{thm:main}): the seed's factor $1/K$ is lost from the deficiency term and retained in
full by the selection term. The first form inside the brace is the direct transcription of the
hidden-seed argument; the second is sensitive to the source's entropy and is the better of the two
whenever $D$ is large or $\epsilon$ is near its minimum $1/D$, where it has no $\log D$ in it at
all. In the closed form of Conjecture~\ref{conj:main} the correct statement is
$\Adv^{\extp} \le c\sqrt{\epsilon R + \log_2 D / K}$, with the $1/K$ attached to the selection term
alone, and $c = 2$ still suffices. Third, Conjecture~\ref{conj:stmt-pub}, the public-seed
conjecture this paper resolves, which already attaches the $1/K$ to the selection term, follows from
Theorem~\ref{thm:main} with $c = 2$ (Corollary~\ref{cor:pub}); so what
Proposition~\ref{prop:false} refutes is the transposition, not that conjecture. Sections and lemmas are kept in one-to-one correspondence with
the hidden-seed note so that the two can be read against each other.

\section{Setting and the two games}\label{sec:setting}

Fix nonempty finite sets $\calK$ (seeds), $\calD$ (inputs) and $\calR$ (outputs), and write
$K := |\calK|$, $D := |\calD|$, $R := |\calR|$. Let $\Fun(\calK \times \calD, \calR)$ be the set of
all functions $\calK \times \calD \to \calR$. For distributions $P, Q$ on a countable set $\Omega$,
\(
  \SD(P,Q) := \max_{A \subseteq \Omega}\bigl(P(A) - Q(A)\bigr)
            = \tfrac12 \sum_{\omega \in \Omega} |P(\omega) - Q(\omega)|
\)
denotes statistical distance, the maximum being attained at
$A^{*} = \{\omega : P(\omega) > Q(\omega)\}$ whether $\Omega$ is finite or not, and $U_{\calR}$ is the uniform distribution on $\calR$. A
\emph{source} $S$ receives a function table $H$ and outputs an input $x \in \calD$ together with an
auxiliary string $z \in \{0,1\}^{*}$. The source $S$, the predictor $P$ and the distinguisher $\mathsf{D}$ of
Figure~\ref{fig:games} are all computationally unbounded, and all three receive the \emph{entire}
table of $H$, all $K \cdot D$ entries, as an explicit input; there is no query budget. The
distinguisher now also receives the seed $\sd$, so the only object hidden from $\mathsf{D}$ is the
coins of $S$.

\subsection{Standard facts}\label{sec:facts}

The seven facts below are the whole of what Part~I borrows from outside itself, apart from the
one result of the companion note quoted as Fact~\ref{fact:hidden}. They are collected here so
that the arguments of Sections~\ref{sec:false} and~\ref{sec:proof} rest on nothing that is not
written down.

\begin{fact}[Optimal distinguishing advantage]\label{fact:sd}
Let $P_0$ and $P_1$ be distributions on a countable set $\Omega$. Sample $b \sample \{0,1\}$ and
then $\omega \sample P_b$. For every function $\mathsf{A} \colon \Omega \to \{0,1\}$, possibly
randomised and subject to no complexity bound,
\[
  2\Pr[\mathsf{A}(\omega) = b] - 1 \;\le\; \SD(P_0,P_1) ,
\]
with equality for the maximum-a-posteriori test, which returns $1$ exactly when
$P_1(\omega) > P_0(\omega)$. In particular the maximum over $\mathsf{A}$ is attained.
\end{fact}

\begin{fact}[Bounded differences; McDiarmid~\cite{McD89}]\label{fact:mcd}
Let $V_1, \dots, V_N$ be independent random variables, $V_i$ taking values in a set
$\mathcal{V}_i$, and let the measurable function
$g \colon \mathcal{V}_1 \times \cdots \times \mathcal{V}_N \to \mathbb{R}$ satisfy
\[
  |g(v) - g(v')| \;\le\; c_i
  \qquad\text{whenever $v$ and $v'$ differ in the $i$th coordinate only.}
\]
Then for every $\lambda > 0$,
\[
  \Pr\bigl[\, g(V_1,\dots,V_N) \;\ge\; \mathbb{E}\,g(V_1,\dots,V_N) + \lambda \,\bigr]
  \;\le\; \exp\Bigl(-\frac{2\lambda^{2}}{\sum_{i=1}^{N} c_i^{2}}\Bigr) .
\]
\end{fact}

\begin{fact}[Binomial estimate]\label{fact:binom}
For integers $1 \le t \le D$, \; $\binom{D}{t} \le (eD/t)^{t}$.
\end{fact}

\begin{fact}[Mean absolute deviation]\label{fact:mad}
For a real random variable $X$ of finite variance,
$\mathbb{E}\,|X - \mathbb{E}X| \le \sqrt{\operatorname{Var} X}$.
\end{fact}

\begin{fact}[Jensen's inequality]\label{fact:jensen}
Let $I \subseteq \mathbb{R}$ be an interval, $\varphi \colon I \to \mathbb{R}$ concave, and $X$
an integrable random variable taking values in $I$. Then
$\mathbb{E}[\varphi(X)] \le \varphi(\mathbb{E}X)$. If $\varphi$ is in addition nondecreasing and
$\mathbb{E}X \le m$ with $m \in I$, then $\mathbb{E}[\varphi(X)] \le \varphi(m)$.
\end{fact}

\begin{fact}[Cauchy--Schwarz]\label{fact:cs}
For reals $\alpha_1,\alpha_2,\beta_1,\beta_2$, \;
$(\alpha_1\beta_1 + \alpha_2\beta_2)^{2} \le (\alpha_1^{2}+\alpha_2^{2})(\beta_1^{2}+\beta_2^{2})$.
\end{fact}

\begin{fact}[Coverage of a random row]\label{fact:cover}
The map $R \mapsto (1-1/R)^{R}$ is increasing on the integers $R \ge 2$, with value $\tfrac14$ at
$R = 2$ and supremum $e^{-1}$. In particular $(1-1/R)^{R} \ge \tfrac14$ for every integer
$R \ge 2$.
\end{fact}


\begin{figure}[ht]
\centering
\fbox{\begin{minipage}[t]{0.36\textwidth}
\underline{$\Gm\ \Pred^{S}_{P}$}\\[2pt]
$H \sample \Fun(\calK \times \calD, \calR)$\\
$(x,z) \sample S(H)$\\
$x' \sample P(H,z)$\\
$\ret (x = x')$
\end{minipage}}\hfill
\fbox{\begin{minipage}[t]{0.55\textwidth}
\underline{$\Gm\ \Extp^{S}_{\mathsf{D}}$}\\[2pt]
$H \sample \Fun(\calK \times \calD, \calR)$\\
$(x,z) \sample S(H)$;\quad $\sd \sample \calK$\\
$y_0 \leftarrow H(\sd,x)$;\quad $y_1 \sample \calR$;\quad $b \sample \{0,1\}$\\
$b' \sample \mathsf{D}(H, \sd, y_b, z)$\\
$\ret (b = b')$
\end{minipage}}
\caption{Prediction game (left) and extraction game (right). The first argument of $S$, $P$ and
$\mathsf{D}$ is the \emph{full function table} of $H$; the seed $\sd$ \emph{is} given to
$\mathsf{D}$, but is sampled only after $S$ has terminated.}
\label{fig:games}
\end{figure}

\begin{definition}\label{def:games}
For the games of Figure~\ref{fig:games}, set
\[
  \Adv^{\pred}_{\calK,\calD,\calR,S}(P) := \Pr[\Pred^{S}_{P} \Rightarrow 1] ,
  \qquad
  \Adv^{\extp}_{\calK,\calD,\calR}(S,\mathsf{D}) := 2\Pr[\Extp^{S}_{\mathsf{D}} \Rightarrow 1] - 1 .
\]
In $\Gm\ \Extp$ the distinguisher is called as $b' \sample \mathsf{D}(H, \sd, y_b, z)$: it is
given the seed. The suffix
$\mathrm{pub}$ is carried throughout to distinguish this from the game $\Gm\ \Ext$ of the
companion notes, which is identical except that $\mathsf{D}$ is called as
$b' \sample \mathsf{D}(H, y_b, z)$ and the seed is withheld; the two advantages
$\Adv^{\extp}$ and $\Adv^{\ext}$ are different quantities and no result about one is a result
about the other. A source $S$ is \emph{$\epsilon$-unpredictable} if
$\Adv^{\pred}_{\calK,\calD,\calR,S}(P) \le \epsilon$ for every unbounded predictor $P$; this is a
property of $S$ alone and is the same notion in both models.
\end{definition}

The hidden-seed game just described is the subject of the companion note, whose bound is quoted
here for the comparison drawn in Remark~\ref{rem:cost}. Nothing in this paper is proved about
$\Adv^{\ext}$, and Fact~\ref{fact:hidden} is used nowhere in the proofs.

\begin{fact}[Companion note, hidden-seed model]\label{fact:hidden}
For all nonempty finite $\calK, \calD, \calR$, all $\epsilon \in (0,1]$, every
$\epsilon$-unpredictable source $S$ and every unbounded distinguisher $\mathsf{D}$ in the game
$\Gm\ \Ext$ of Definition~\ref{def:games},
\[
  \Adv^{\ext}_{\calK,\calD,\calR}(S,\mathsf{D})
  \;\le\; \frac{1}{\sqrt{2}}\sqrt{\frac{\epsilon R}{K}}
       \;+\; \frac{6}{5}\sqrt{\frac{1 + \ln D}{K}} .
\]
\end{fact}

The prediction game is untouched by the change, and the choice not to hand $\sd$ to $P$ costs
nothing: $\sd$ is sampled after $S$ terminates and is therefore independent of $(X,H,Z)$, so a
predictor given $\sd$ has the same optimal success probability. What the change does affect is
the sampling order in $\Gm\ \Extp$, and that order is now the entire content of the model. The
source must commit to $x$ before $\sd$ exists; were $S$ allowed to see $\sd$ it could choose $x$
with $H(\sd,x)$ whatever value it liked and no bound of any shape would hold.

Write $(X,Z)$ for the output distribution of $S$, and for a view $(H,z)$ of positive probability put
\[
  p^{H,z}_{x} \;:=\; \Pr[X = x \mid H, z] ,
  \qquad
  \epsilon_{H,z} \;:=\; \max_{x \in \calD} p^{H,z}_{x} .
\]
For $k \in \calK$ let $\nu^{H,z}_{k}$ denote the conditional law of $H(k, X)$ given $(H,Z) = (H,z)$,
that is, the pushforward of $p^{H,z}$ along the single row $H(k,\cdot)$, and put
\[
  \Delta_{H,z} \;:=\; \frac{1}{K}\sum_{k \in \calK} \SD\bigl(\nu^{H,z}_{k}, U_{\calR}\bigr).
\]

\begin{lemma}[Analytic forms]\label{lem:analytic}
It holds that
$\max_{P} \Adv^{\pred}_{\calK,\calD,\calR,S}(P) = \mathbb{E}_{(H,Z)}[\epsilon_{H,Z}]$, so that
$S$ is $\epsilon$-unpredictable if and only if $\mathbb{E}_{(H,Z)}[\epsilon_{H,Z}] \le \epsilon$;
and
\begin{equation}\label{eq:analytic}
  \max_{\mathsf{D}}\Adv^{\extp}_{\calK,\calD,\calR}(S,\mathsf{D})
  \;=\; \mathbb{E}_{(H,Z)}\bigl[\Delta_{H,Z}\bigr] .
\end{equation}
\end{lemma}

\begin{proof}
% BEGIN overview lem:analytic
Each claim identifies an optimal unbounded adversary with a statistical quantity. For the first,
the predictor that outputs a mode of $p^{H,z}$ succeeds with probability $\epsilon_{H,z}$ on the
view $(H,z)$, and its success is averaged over $(H,Z)$. For the second, the argument conditions on
$(H,Z) = (H,z)$, identifies the optimal advantage with the statistical distance between the two
view distributions, and expands that distance over the pair $(\sd,y_b)$; since $\sd$ is uniform
under both hypotheses, the sum factors through the $K$ rows and equals $\Delta_{H,z}$. Averaging
over $(H,Z)$ gives \eqref{eq:analytic}.
% END overview lem:analytic

The first claim is unchanged from the hidden-seed setting: given its view $(H,z)$ a predictor
succeeds with probability at most $\epsilon_{H,z}$, with equality for the (unbounded) predictor
that outputs a mode of $p^{H,z}$; averaging over $(H,Z)$ gives the claim. (The view space
$\Fun(\calK \times \calD, \calR) \times \{0,1\}^{*}$ is countably infinite, but the maxima over
$P$, and below over $\mathsf{D}$, are attained all the same: the mode predictor and the
maximum-a-posteriori test of Fact~\ref{fact:sd} are defined view by view.)

For the second, condition on $(H,Z) = (H,z)$. The view of $\mathsf{D}$ is $(H, \sd, y_b, z)$, and
the two hypotheses differ only in the pair $(\sd, y_b)$: under $b = 0$ it is distributed as
$(\sd, H(\sd,X))$ with $\sd \sample \calK$ independent of $X$, and under $b = 1$ as
$U_{\calK} \otimes U_{\calR}$. By Fact~\ref{fact:sd} the optimal advantage $2\Pr[b'=b]-1$ equals the
statistical distance between the two view distributions; since the component $(H,Z)$ has the same
law under $b=0$ and $b=1$, and the first coordinate $\sd$ is uniform under both hypotheses,
\begin{align*}
  \SD\bigl((\sd, H(\sd,X)) \mid (H,z),\; U_{\calK}\otimes U_{\calR}\bigr)
  &= \tfrac12 \sum_{k,r} \tfrac1K \bigl|\nu^{H,z}_{k}(r) - \tfrac1R\bigr| \\
  &= \frac1K \sum_{k} \SD\bigl(\nu^{H,z}_{k}, U_{\calR}\bigr)
  \;=\; \Delta_{H,z},
\end{align*}
and averaging over $(H,Z)$ gives \eqref{eq:analytic}.
\end{proof}

The one structural change is visible in the display above: the average over $k$ now sits
\emph{outside} the statistical distance. In the hidden-seed model $\mathsf{D}$ sees a mixture of
the $K$ rows and the mixing itself smooths the distribution; here $\mathsf{D}$ is told which row it
is looking at, and the $K$ rows are graded separately. Everything below follows from that one
displacement.

\section{The conjectured bound is false in this model}\label{sec:false}

\begin{proposition}\label{prop:false}
For every $c > 0$ there are nonempty finite $\calK, \calD, \calR$, a value $\epsilon \in (0,1]$, an
$\epsilon$-unpredictable source $S$ and a distinguisher $\mathsf{D}$ with
\[
  \Adv^{\extp}_{\calK,\calD,\calR}(S,\mathsf{D}) \;\ge\; \tfrac14
  \;>\; c\,\sqrt{\frac{\epsilon R + \log_2 D}{K}} .
\]
The witness has $z = \bot$ and $X$ uniform on all of $\calD$, that is, the largest entropy a
source over $\calD$ can have.
\end{proposition}

\begin{proof}
% BEGIN overview prop:false
A witness is exhibited and its two sides are computed. The source outputs a uniform $x$ and
$z = \bot$, and $\calR$ is taken with $R = D$; by Lemma~\ref{lem:analytic} it is
$(1/D)$-unpredictable, so $\epsilon R = 1$ and the conjectured bound reads
$c\sqrt{(1+\log_2 D)/K}$. The distinguisher returns $0$ when its challenge lies in the image of the
row $H(\sd,\cdot)$, which its view determines. Under $b = 0$ this is correct with probability $1$;
under $b = 1$ it is correct with probability $(1-1/R)^{D}$, computed from the $D$ i.i.d.\ entries
of the row and bounded below by Fact~\ref{fact:cover}. Combining the two gives advantage at least $\tfrac14$ for every $K$, and $K$ is then
taken large enough to put the conjectured bound below $\tfrac14$.
% END overview prop:false

Take $\calD$ and $\calR$ with $R = D \ge 2$, let $S(H)$ output $x \sample \calD$ and
$z \leftarrow \bot$, and let $\calK$ be arbitrary. For each $(H,z)$ the output $X$ is uniform on
$\calD$, so $\epsilon_{H,z} = 1/D$ and $S$ is $(1/D)$-unpredictable by
Lemma~\ref{lem:analytic}; thus $\epsilon = 1/D$ and $\epsilon R = 1$. Let $\mathsf{D}(H,\sd,y,z)$
return $0$ if $y$ lies in the image of the row $H(\sd,\cdot)$ and $1$ otherwise. The test is
computable from the view because $\mathsf{D}$ holds both $H$ and $\sd$. If $b = 0$ then
$y_0 = H(\sd,X)$ is in that image by construction, so $\mathsf{D}$ is correct with probability $1$.
If $b = 1$ then $y_1 \sample \calR$ is independent of $H$ and $\sd$, and $\mathsf{D}$ is correct
with probability $\mathbb{E}\bigl[1 - |\mathrm{Im}\,H(\sd,\cdot)|/R\bigr]$; the $D$ entries of the
row are i.i.d.\ uniform on $\calR$, so a fixed $r$ is missed with probability $(1-1/R)^{D}$ and this
expectation equals $(1-1/R)^{D} = (1-1/R)^{R}$, which is at least $\tfrac14$ by
Fact~\ref{fact:cover}. Hence
$\Adv^{\extp} = 2 \cdot \tfrac12\bigl(1 + (1-1/R)^{R}\bigr) - 1 \ge \tfrac14$, uniformly in $K$.
Meanwhile $c\sqrt{(\epsilon R + \log_2 D)/K} = c\sqrt{(1+\log_2 D)/K} < \tfrac14$ as soon as
$K > 16c^{2}(1 + \log_2 D)$, and $K$ is unconstrained.
\end{proof}

The constant $\tfrac14$ is attained rather than approached: by Fact~\ref{fact:cover} the
proposition in fact gives
$\Adv^{\extp} = (1-1/R)^{R} \ge \tfrac14$.

The source of the failure is not the auxiliary output and not adversarial selection: it is that
once $\sd$ is public, the map $x \mapsto H(\sd,x)$ is known to $\mathsf{D}$ in full, so the only
randomness available is that of $X$ itself and the seed cannot contribute entropy. Any correct
bound must therefore contain a term in $\epsilon R$ that does not decrease in $K$. That is the
only defect: with the $1/K$ moved onto the selection term alone, Conjecture~\ref{conj:main} holds
with $c = 2$.

\section{The corrected bound}\label{sec:corrected}

\begin{theorem}\label{thm:main}
For all nonempty finite $\calK, \calD, \calR$, all $\epsilon \in (0,1]$, every $\epsilon$-unpredictable
source $S$ and every unbounded distinguisher $\mathsf{D}$ in the game of Figure~\ref{fig:games},
\begin{align}
  \Adv^{\extp}_{\calK,\calD,\calR}(S,\mathsf{D})
  &\;\le\; \frac{1}{\sqrt{2}}\sqrt{\epsilon R}
        \;+\; \frac{6}{5}\sqrt{\frac{1 + \ln D}{K}}
  \;\le\; 2\sqrt{\epsilon R + \frac{\log_2 D}{K}} \; ,
  \label{eq:main-plain}\\[2pt]
  \Adv^{\extp}_{\calK,\calD,\calR}(S,\mathsf{D})
  &\;\le\; \frac{1}{\sqrt{2}}\sqrt{\epsilon R}
        \;+\; \sqrt{\frac{\ln(2eD\epsilon)}{2K}} \;+\; \frac{9}{10\sqrt{K}} \; .
  \label{eq:main-sharp}
\end{align}
As in the hidden-seed setting the last inequality of \eqref{eq:main-plain} is not an unconditional
fact about the parameters: it holds for all $\epsilon \in (0,1]$ whenever $D \ge 2$, and for
$D = 1$ it holds at $\epsilon = 1$, the only value at which an $\epsilon$-unpredictable source
exists; it can fail for $D = 1$ and $\epsilon$ small. Neither of
\eqref{eq:main-plain} and \eqref{eq:main-sharp} implies the other; see Remark~\ref{rem:refine}.
\end{theorem}

\begin{theoremstar}[Informal; see Theorem~\ref{thm:main}]
Publishing the seed costs exactly one factor of $K$, and costs it in one summand only. The first
summand $\tfrac{1}{\sqrt{2}}\sqrt{\epsilon R}$ does not improve as the seed space grows; the
second, which is what the source can extract by choosing its support after seeing $H$ and signalling
it through $z$, keeps its $1/K$ in full. The two forms differ only in that second summand:
\eqref{eq:main-plain} carries $1 + \ln D$ where \eqref{eq:main-sharp} carries the
entropy-sensitive $\ln(2eD\epsilon)$, and \eqref{eq:main-sharp} pays an additive
$\tfrac{9}{10\sqrt{K}}$ for the exchange.
\end{theoremstar}

Section~\ref{sec:posed} reproduces two conjectures, one for each game, and neither of them is
refuted here: what Proposition~\ref{prop:false} kills is Conjecture~\ref{conj:main}, the
transposition of the secret-seed Conjecture~\ref{conj:stmt-secret} to the public-seed game. The
public-seed conjecture as actually posed is the weaker Conjecture~\ref{conj:stmt-pub},
\begin{equation}\label{eq:conj}
  \Adv^{\extp}_{\calK,\calD,\calR}(S,\mathsf{D})
  \;\le\; c\Bigl(\sqrt{\epsilon R} \;+\; \sqrt{\tfrac{\log_2 D}{K}}\Bigr) ,
\end{equation}
which already carries the factor $1/K$ on the selection term alone. That conjecture survives, and
Theorem~\ref{thm:main} settles it.

\begin{corollary}[The public-seed conjecture as posed]\label{cor:pub}
For all nonempty finite $\calK, \calD, \calR$ with $D \ge 2$, all $\epsilon \in (0,1]$, every
$\epsilon$-unpredictable source $S$ and every unbounded distinguisher $\mathsf{D}$, the bound
\eqref{eq:conj} holds with $c = \tfrac85$; the smallest constant obtainable by this route is
$\tfrac65\sqrt{1 + \ln 2} = 1.5615\ldots$, attained at $D = 2$. For $D = 1$, the only case left,
\eqref{eq:conj} holds with $c = \tfrac{1}{\sqrt2} + \tfrac65 < 1.91$. Hence $c = 2$ suffices
throughout the range in which an $\epsilon$-unpredictable source exists.
\end{corollary}

\begin{proof}
% BEGIN overview cor:pub
For $D \ge 2$ the two summands of \eqref{eq:main-plain} are compared with the two summands of
\eqref{eq:conj} separately: the first comparison asks for $c \ge 1/\sqrt2$, and the second
reduces to bounding $(1+\ln D)/\log_2 D$, which is decreasing in $D$ and so largest at $D = 2$.
For $D = 1$ the second summand of \eqref{eq:conj} vanishes, so the comparison is made against
$\sqrt{\epsilon R}$ alone, using that $D = 1$ forces $\epsilon = 1$ and hence $\epsilon R \ge 1$.
% END overview cor:pub

By \eqref{eq:main-plain} it suffices that $\tfrac{1}{\sqrt2} \le c$ and that
$\tfrac65\sqrt{(1+\ln D)/K} \le c\sqrt{\log_2 D / K}$ for every $D \ge 2$, that is, that
$c \ge \tfrac65\sqrt{h(D)}$ where $h(D) := (1+\ln D)/\log_2 D$. Writing
$\log_2 D = \ln D/\ln 2$ gives $h(D) = (\ln 2)(1 + 1/\ln D)$, which is decreasing on $D \ge 2$;
so $h(D) \le h(2) = (\ln 2)(1 + 1/\ln 2) = \ln 2 + 1 = 1.6931\ldots$ and
$\tfrac65\sqrt{h(2)} = 1.5615\ldots$. Since $1.5615 \ge 1/\sqrt2$, both requirements hold at
$c = 1.5615\ldots$, and a fortiori at $c = \tfrac85$.

For $D = 1$ the second summand of \eqref{eq:conj} vanishes, so what is needed is
$\tfrac{1}{\sqrt2}\sqrt{\epsilon R} + \tfrac65 K^{-1/2} \le c\sqrt{\epsilon R}$. Here $\epsilon = 1$,
because $\mathbb{E}[\epsilon_{H,Z}] \ge 1/D = 1$ leaves no smaller value admissible, so
$\epsilon R = R \ge 1$; together with $K \ge 1$ this gives
$\tfrac{1}{\sqrt2}\sqrt{\epsilon R} + \tfrac65 K^{-1/2} \le
\bigl(\tfrac{1}{\sqrt2} + \tfrac65\bigr)\sqrt{\epsilon R} < 1.91\sqrt{\epsilon R}$.
\end{proof}

Proposition~\ref{prop:false} therefore refutes the verbatim transposition of the secret-seed
conjecture, which is strictly stronger than \eqref{eq:conj}, since
$\sqrt{(\epsilon R + \log_2 D)/K} \le \sqrt{\epsilon R} + \sqrt{\log_2 D/K}$. Both are recorded
because the transposition is the natural first guess, and locating where it fails is what
identifies the term that cannot decay in $K$.

The rest of Part~I proves Theorem~\ref{thm:main}: Section~\ref{sec:setup} fixes the notation,
Section~\ref{sec:lemmas} establishes the three lemmas the proof rests on, and
Section~\ref{sec:assemble} puts them together. Section~\ref{sec:disc} discusses what the public
seed does and does not cost, and Part~II asks what the bound so obtained buys.

\section{Proof of Theorem~\ref{thm:main}}\label{sec:proof}

\subsection{Set-up}\label{sec:setup}

For a function table $H$, a seed $k \in \calK$ and a nonempty $T \subseteq \calD$ let
$\emp^{H}_{k,T}$ be the empirical distribution of the $|T|$ oracle values in row $k$ over $T$,
\[
  \emp^{H}_{k,T}(r) \;:=\; \frac{1}{|T|}\,\bigl|\{x \in T : H(k,x) = r\}\bigr| ,
\]
and put
\[
  F_{T}(H) \;:=\; \frac{1}{K}\sum_{k \in \calK} \SD\bigl(\emp^{H}_{k,T}, U_{\calR}\bigr),
  \qquad
  \Phi_{H}(t) \;:=\; \max_{T \subseteq \calD,\; |T| = t} F_{T}(H) .
\]
For a uniform $H$ and each fixed $k$ the $|T|$ values $\{H(k,x)\}_{x \in T}$ are $|T|$ i.i.d.\
uniform draws from $\calR$; the $K$ rows are independent of one another, but that is not needed for
the mean, only for the deviation. The pair $(F_{T}, \Phi_{H})$ is the exact point of departure from the hidden-seed proof,
where the corresponding object $Q_{H,T}$ was a single empirical distribution of $Kt$ draws. Here
there are $K$ separate empirical distributions of $t$ draws each, and their distances are averaged.

\smallskip
The proof is four steps, in the same order as before: flatten the source (Lemma~\ref{lem:flat}),
compute the mean for one fixed support (Lemma~\ref{lem:mean}), control the deviation \emph{uniformly
over all supports of all sizes} (Lemma~\ref{lem:unif}), and average (Section~\ref{sec:assemble}).
Lemma~\ref{lem:flat} and part~(i) of Lemma~\ref{lem:unif} carry over from the hidden-seed note with
$F_{\cdot}(H)$ in place of $\SD(Q_{H,\cdot},U_{\calR})$ and the centering of
Lemma~\ref{lem:mean} in place of the old one (the deviation argument is untouched, but the
centering is inherited, so $\tfrac12\sqrt{R/(Kt)}$ becomes $\tfrac12\sqrt{R/t}$ inside $W_{1}$),
and the convexity step of Lemma~\ref{lem:flat} now has to pass through the average over $k$. Of the
substance, only Lemma~\ref{lem:mean} changes, and it changes by exactly one factor of $K$.
Part~(ii) of Lemma~\ref{lem:unif} is new and applies verbatim to the hidden-seed proof as well.

\subsection{The three lemmas}\label{sec:lemmas}

\begin{lemma}[Flattening]\label{lem:flat}
Fix $H$ and $z$, and let $t := \lfloor 1/\epsilon_{H,z} \rfloor$. Then $1 \le t \le D$ and
\(
  \Delta_{H,z} \le \Phi_{H}(t).
\)
\end{lemma}

The effect is that the source enters the rest of the argument only through the single integer $t$:
Lemmas~\ref{lem:mean} and~\ref{lem:unif} are statements about flat supports $T$ with $|T| = t$ and
about $H$, and the unpredictability hypothesis reappears only when $t$ is averaged over $(H,Z)$ in
Section~\ref{sec:assemble}.

\begin{proof}
% BEGIN overview lem:flat
The floor gives $1 \le t \le D$ and the pointwise cap $\|p^{H,z}\|_{\infty} \le 1/t$. The
distributions obeying that cap form a hypersimplex, whose vertices are identified as the uniform
distributions $\tfrac1t\mathbf{1}_{T}$ on $t$-element supports, using that $t$ is an integer; so
$p^{H,z}$ is a convex combination of them. The map $p \mapsto \Delta_{H,z}$ is then shown to be
convex, being an average over $k$ of a convex function of a linear pushforward, and its value at
the vertex $\tfrac1t\mathbf{1}_{T}$ is $F_{T}(H)$. Convexity along the combination bounds
$\Delta_{H,z}$ by $\Phi_{H}(t)$.
% END overview lem:flat

Since $\epsilon_{H,z} = \max_x p^{H,z}_x \in [1/D, 1]$ we have $1 \le t \le D$, and $\|p^{H,z}\|_\infty
\le 1/t$ by definition of the floor. The set $\{p \in \mathbb{R}^{\calD} : p \ge 0,\ \sum_x p_x = 1,\
p_x \le 1/t\}$ is the hypersimplex, whose vertices are precisely the vectors $\tfrac{1}{t}\mathbf{1}_T$
with $|T| = t$: a vertex has at most one coordinate strictly inside $(0, 1/t)$, and since $t$ is an
integer the mass constraint rules that coordinate out. (Integrality of $t$ is essential here; a
polytope with a non-integral cap has fractional vertices.) Hence $p^{H,z} = \sum_j \alpha_j
\tfrac{1}{t}\mathbf{1}_{T_j}$ for some convex weights $\alpha_j$ and $t$-sets $T_j$. For each fixed
$k$ the map $p \mapsto \nu_{k}$ is linear, being a pushforward, and $\SD(\cdot,U_{\calR})$ is convex;
a nonnegative average over $k$ of convex functions of linear maps is convex, so $p \mapsto
\Delta_{H,z}$ is convex in $p^{H,z}$. Finally the value of that map at the vertex
$\tfrac1t\mathbf{1}_{T}$ is $F_{T}(H)$, since the pushforward of $\tfrac1t\mathbf{1}_{T}$ along
$H(k,\cdot)$ is $\emp^{H}_{k,T}$. Hence
$\Delta_{H,z} \le \sum_j \alpha_j F_{T_j}(H) \le \Phi_H(t)$.
\end{proof}

\begin{lemma}[Mean at a fixed support]\label{lem:mean}
For every fixed $T$ with $|T| = t$, \;
$\mathbb{E}_{H}\bigl[F_{T}(H)\bigr] \le \tfrac12 \sqrt{R/t}$.
\end{lemma}

\begin{proof}
% BEGIN overview lem:mean
The distance is bounded one row at a time. For fixed $k$, each $\emp^{H}_{k,T}(r)$ is an average of
$t$ i.i.d.\ Bernoulli$(1/R)$ variables, so its mean absolute deviation from $1/R$ is bounded by its
standard deviation; summing over the $R$ values of $r$ and halving gives $\tfrac12\sqrt{R/t}$ for
that row, and the average over the $K$ seeds inherits the bound.
% END overview lem:mean

Fix $k$. Each $\emp^{H}_{k,T}(r)$ is an average of $t$ i.i.d.\ Bernoulli$(1/R)$
variables, so Fact~\ref{fact:mad} gives $\mathbb{E}\,|\emp^{H}_{k,T}(r) - 1/R| \le \sqrt{\operatorname{Var}} =
\sqrt{\tfrac{1}{t}\tfrac1R(1-\tfrac1R)} \le (tR)^{-1/2}$. Summing over the $R$ values of $r$ and
halving gives $\mathbb{E}\,\SD(\emp^{H}_{k,T},U_{\calR}) \le \tfrac12 R (tR)^{-1/2} =
\tfrac12\sqrt{R/t}$, and averaging over the $K$ seeds preserves the bound.
\end{proof}

Lemma~\ref{lem:mean} is the whole of the loss. Its hidden-seed version averages $Kt$ draws and
yields $\tfrac12\sqrt{R/(Kt)}$; here the count is $t$, because each row is graded on its own $t$
values and the average over $k$ happens afterwards, outside the distance.

\begin{lemma}[Uniform deviation]\label{lem:unif}
Let $a := 1 + \ln D$, let $b_{t} := \tfrac1t \ln\binom{D}{t}$ for $1 \le t \le D$, put
$\psi(u) := e^{-u}/(1-e^{-u})$, and set
\[
  W_{1} \;:=\; \max_{1 \le t \le D}\Bigl(\Phi_{H}(t) - \tfrac12\sqrt{\tfrac{R}{t}}\Bigr),
  \qquad
  W_{2} \;:=\; \max_{1 \le t \le D}\Bigl(\Phi_{H}(t) - \tfrac12\sqrt{\tfrac{R}{t}}
                                          - \sqrt{\tfrac{b_{t}}{2K}}\Bigr),
\]
with $W_{i,+} := \max\{W_{i}, 0\}$. Then for every $u > 0$,
\[
  \text{(i)}\quad \Pr_{H}\Bigl[\, W_{1} > \sqrt{\tfrac{a+u}{2K}} \,\Bigr] \;\le\; \psi(u) ,
  \qquad\qquad
  \text{(ii)}\quad \Pr_{H}\Bigl[\, W_{2} > \sqrt{\tfrac{u}{2K}} \,\Bigr] \;\le\; \psi(u) ,
\]
and hence
\[
  \text{(i)}\quad
  \mathbb{E}_{H}[W_{1,+}] \;\le\; \sqrt{\frac{a}{2K}} + \frac{\ln 4}{2\sqrt{2Ka}}
  \;\le\; \frac{6}{5}\sqrt{\frac{a}{K}} ,
  \qquad
  \text{(ii)}\quad
  \mathbb{E}_{H}[W_{2,+}] \;\le\; \frac{3}{2}\sqrt{\frac{\ln 2}{2K}} \;\le\; \frac{9}{10\sqrt{K}} .
\]
\end{lemma}

The two parts bound the same quantity twice and differ only in how the union over supports is paid
for. Part~(i) subtracts from $\Phi_{H}(t)$ a single cost $a = 1 + \ln D$ shared by every size $t$;
part~(ii) subtracts the size-dependent $b_{t} = \tfrac1t\ln\binom{D}{t}$, which is what
$\binom{D}{t}$ costs at that size and which is $0$ at $t = D$. Part~(i) is what yields
\eqref{eq:main-plain} and part~(ii) is what yields \eqref{eq:main-sharp}.

\begin{proof}
% BEGIN overview lem:unif
The argument runs in two phases. First, for a fixed $T$ of size $t$, $F_{T}(H)$ is regarded as a
function of the $Kt$ table entries it depends on; a single entry changes it by at most $1/(Kt)$, so
Fact~\ref{fact:mcd} bounds $\Pr[F_{T}(H) \ge \mathbb{E}\,F_{T}(H) + \lambda]$ by
$e^{-2\lambda^{2}Kt}$, and Lemma~\ref{lem:mean} replaces the mean by $\tfrac12\sqrt{R/t}$.

Second, that bound is summed over all $T$ of all sizes under two choices of $\lambda$. Part~(i)
takes $\lambda = \sqrt{(a+u)/(2K)}$, which does not depend on $t$, and absorbs $\binom{D}{t}$ into
$e^{ta}$; part~(ii) takes $\lambda_{t} = \sqrt{(b_{t}+u)/(2K)}$, for which that bound
becomes $\binom{D}{t}^{-1}e^{-tu}$ and the binomial coefficient cancels against the union bound.
Both leave the geometric sum $\sum_{t \ge 1}e^{-tu} = \psi(u)$. The two expectations then follow by
integrating $\Pr[W_{i,+} > s]$ over $s$, substituting for $u$ and splitting each integral at
$\ln 2$, the point where $\psi(u) = 1$.
% END overview lem:unif

Fix $T$ with $|T| = t$ and regard $F_{T}(H)$ as a function of the $Kt$ independent values
$\{H(k,x)\}_{k \in \calK, x \in T}$. Changing one of them, say $H(k_0,x_0)$, from $r_1$ to $r_2$
affects only the summand $k = k_0$: it shifts $\emp^{H}_{k_0,T}(r_1)$ by $-1/t$ and
$\emp^{H}_{k_0,T}(r_2)$ by $+1/t$ and leaves the other coordinates alone, so it changes
$\tfrac12\sum_r|\emp^{H}_{k_0,T}(r) - 1/R|$ by at most $1/t$, and changes $F_T(H)$, which carries
that summand at weight $1/K$, by at most $1/(Kt)$. The bounded-differences constant is therefore
$1/(Kt)$, exactly as in the hidden-seed setting, and Fact~\ref{fact:mcd} with
$\sum_i c_i^2 = Kt \cdot (Kt)^{-2} = 1/(Kt)$ gives, for every $\lambda > 0$,
\[
  \Pr\bigl[F_{T}(H) \ge \mathbb{E}\,F_{T}(H) + \lambda\bigr]
  \;\le\; \exp\bigl(-2\lambda^{2} K t\bigr) .
\]
The independence hypothesis of Fact~\ref{fact:mcd} is met because $H$ is a uniform \emph{function} table, so its
$KD$ entries are i.i.d.; were each row $H(k,\cdot)$ a random permutation or injection, the $t$
values in a row would be sampled without replacement and the inequality would not apply as stated.
The same proviso governs the hidden-seed note.
Since Lemma~\ref{lem:mean} gives $\mathbb{E}\,F_{T}(H) \le \tfrac12\sqrt{R/t}$, the event
$F_{T}(H) > \tfrac12\sqrt{R/t} + \lambda$ is contained in the event bounded above, for every
$\lambda > 0$ we care to choose. The two parts differ only in that choice.

\emph{(i).} Take $\lambda := \sqrt{(a+u)/(2K)}$, which does \emph{not} depend on $t$; then the
right-hand side is $e^{-t(a+u)}$. Union bounding over all $T$ of all sizes, using
Fact~\ref{fact:binom} in the form $\binom{D}{t} \le (eD/t)^{t} \le e^{t(1+\ln D)} = e^{ta}$,
\[
  \Pr\Bigl[W_{1} > \lambda\Bigr]
  \;\le\; \sum_{t=1}^{D} \binom{D}{t} e^{-t(a+u)}
  \;\le\; \sum_{t \ge 1} e^{ta} e^{-t(a+u)}
  \;=\; \sum_{t\ge1} e^{-tu}
  \;=\; \psi(u).
\]
For the expectation, substitute $s = \sqrt{(a+u)/(2K)}$ in $\mathbb{E}[W_{1,+}] = \int_0^\infty \Pr[W_{1,+} > s]\,ds$,
the range $s \le \sqrt{a/(2K)}$ being bounded by $\Pr \le 1$:
\begin{align*}
  \mathbb{E}[W_{1,+}] \;&\le\; \sqrt{\tfrac{a}{2K}}
  + \int_0^\infty \min\{1,\psi(u)\}\,\frac{du}{2\sqrt{2K(a+u)}} \\
  &\le\; \sqrt{\tfrac{a}{2K}} + \frac{1}{2\sqrt{2Ka}}\int_0^\infty \min\{1,\psi(u)\}\,du .
\end{align*}
Since $\psi(u) \ge 1$ exactly for $u \le \ln 2$, the last integral equals
$\ln 2 + \bigl[\ln(1-e^{-u})\bigr]_{\ln 2}^{\infty} = 2\ln 2 = \ln 4$. Finally $a \ge 1$ gives
$\sqrt{a/(2K)}\,(1 + \tfrac{\ln 4}{2a}) \le \tfrac{1.6932}{\sqrt2}\sqrt{a/K} \le \tfrac65\sqrt{a/K}$.

\emph{(ii).} The bound $\binom{D}{t} \le e^{ta}$ is lossy for large $t$, and charging each size its
true cost costs nothing. Take instead $\lambda_{t} := \sqrt{(b_{t}+u)/(2K)}$, which does depend on
$t$; the right-hand side of the bound of Fact~\ref{fact:mcd} is now
$e^{-t(b_{t}+u)} = \binom{D}{t}^{-1}e^{-tu}$, so the binomial coefficient cancels against the union
bound instead of being dominated by it:
\[
  \Pr\Bigl[W_{2} > \sqrt{\tfrac{u}{2K}}\Bigr]
  \;\le\; \sum_{t=1}^{D}\binom{D}{t}\binom{D}{t}^{-1}e^{-tu}
  \;\le\; \sum_{t\ge1} e^{-tu} \;=\; \psi(u) ,
\]
where the first inequality uses that $W_{2} > \sqrt{u/(2K)}$ forces
$\Phi_{H}(t) - \tfrac12\sqrt{R/t} > \sqrt{b_{t}/(2K)} + \sqrt{u/(2K)} \ge \lambda_{t}$ for some
$t$, by subadditivity of the square root. For the expectation, substitute $s = \sqrt{u/(2K)}$, so
that $ds = du/(2\sqrt{2Ku})$ and no range needs separate treatment:
\[
  \mathbb{E}[W_{2,+}] \;\le\; \int_{0}^{\infty}\min\{1,\psi(u)\}\,\frac{du}{2\sqrt{2Ku}}
  \;=\; \frac{1}{2\sqrt{2K}}\int_{0}^{\infty}\frac{\min\{1,\psi(u)\}}{\sqrt u}\,du .
\]
Splitting at $\ln 2$ again, the first piece is $\int_{0}^{\ln 2}u^{-1/2}du = 2\sqrt{\ln 2}$ and the
second is at most $(\ln 2)^{-1/2}\int_{\ln 2}^{\infty}\psi(u)\,du = (\ln 2)^{-1/2}\ln 2 =
\sqrt{\ln 2}$, so the integral is at most $3\sqrt{\ln 2}$ and
$\mathbb{E}[W_{2,+}] \le \tfrac32\sqrt{\ln 2/(2K)} \le \tfrac{9}{10}K^{-1/2}$, the last step
because $\tfrac32\sqrt{\ln 2/2} < 0.8831$.
\end{proof}

Note that $b_{D} = 0$: at $t = D$ part~(ii) charges no selection cost whatever. The assembling step
below nevertheless reintroduces a constant there, through the two estimates
$b_{t} \le \ln(eD/t)$ and $t \ge 1/(2\epsilon_{H,z})$; carrying $b_{t}$ itself further would let
the selection term vanish at $\epsilon = 1/D$, as Remark~\ref{rem:refine} argues it should.

The point to record is that the concentration is \emph{identical} to the hidden-seed case even
though the mean is larger by $\sqrt K$. A single table entry moves $Q_{H,T}$, an average of $Kt$
values, by $1/(Kt)$; it moves $F_T$, an average of $K$ statistics of $t$ values each, by
$\tfrac1K \cdot \tfrac1t$, which is the same number. Revealing the seed costs the mean and leaves
the fluctuation alone.

\subsection{Assembling}\label{sec:assemble}

% BEGIN overview thm:main
Lemma~\ref{lem:flat} and the definitions of $W_{1}$ and $W_{2}$ bound $\Delta_{H,z}$ by
$\tfrac12\sqrt{R/t}$ plus a deviation, in two forms. The estimate $t \ge 1/(2\epsilon_{H,z})$ turns
the first summand into $\tfrac12\sqrt{2R\epsilon_{H,z}}$, and the deviations depend on $H$ alone,
so they average to the constants of Lemma~\ref{lem:unif}. Taking $\mathbb{E}_{(H,Z)}$ and applying
Fact~\ref{fact:jensen} with $\mathbb{E}[\epsilon_{H,Z}] \le \epsilon$ gives \eqref{eq:main-plain};
for \eqref{eq:main-sharp}, $b_{t}$ is first bounded by $\ln(2eD\epsilon_{H,z})$ and
Fact~\ref{fact:jensen} applied a second time, on the range where $s \mapsto \sqrt{\ln(2eDs)}$ is
concave. The closed form of
\eqref{eq:main-plain} follows by bounding the ratio of the two bounds, by Fact~\ref{fact:cs} for
$D \ge 2$ and directly for $D = 1$.
% END overview thm:main

Write $t := \lfloor 1/\epsilon_{H,z}\rfloor$. By Lemma~\ref{lem:flat} and the definitions of
$W_{1}$ and $W_{2}$, for every $(H,z)$,
\[
  \Delta_{H,z} \;\le\; \tfrac12\sqrt{\tfrac{R}{t}} + W_{1,+} ,
  \qquad\text{and}\qquad
  \Delta_{H,z} \;\le\; \tfrac12\sqrt{\tfrac{R}{t}} + \sqrt{\tfrac{b_{t}}{2K}} + W_{2,+} .
\]
For $y \ge 1$ one has $\lfloor y \rfloor \ge y/2$, so $t \ge 1/(2\epsilon_{H,z})$ and the first
summand of each is at most $\tfrac12\sqrt{2 R\,\epsilon_{H,z}}$. Both $W_{1,+}$ and $W_{2,+}$ are
functions of $H$ alone, and the $H$-marginal of the joint $(H,Z)$ is uniform, $Z$ being computed
from $H$ and the source's coins, so $\mathbb{E}_{(H,Z)}[W_{i,+}] = \mathbb{E}_{H}[W_{i,+}]$. Taking
$\mathbb{E}_{(H,Z)}$ in the first inequality, using \eqref{eq:analytic} and Fact~\ref{fact:jensen}
for the concave nondecreasing $\sqrt{\cdot}$ with $\mathbb{E}[\epsilon_{H,Z}] \le \epsilon$,
\[
  \Adv^{\extp}_{\calK,\calD,\calR}(S,\mathsf{D})
  \;\le\; \frac{1}{\sqrt2}\sqrt{\epsilon R} \;+\; \mathbb{E}_H[W_{1,+}]
  \;\le\; \frac{1}{\sqrt2}\sqrt{\epsilon R} + \frac{6}{5}\sqrt{\frac{1+\ln D}{K}} ,
\]
which is the first inequality of \eqref{eq:main-plain}. For \eqref{eq:main-sharp}, use the second:
$b_{t} \le \ln(eD/t)$ by Fact~\ref{fact:binom}, and $t \ge 1/(2\epsilon_{H,z})$ gives
$eD/t \le 2eD\epsilon_{H,z}$, so the middle summand is at most
$\sqrt{\ln(2eD\epsilon_{H,z})/(2K)}$. The function $g(s) := \sqrt{\ln(2eDs)}$ is concave wherever
$\ln(2eDs) > 0$, and $\epsilon_{H,z} \ge 1/D$ gives $2eD\epsilon_{H,z} \ge 2e$, so Fact~\ref{fact:jensen}
applies on the range of $\epsilon_{H,Z}$; being also nondecreasing there, $g$ satisfies
$\mathbb{E}\,g(\epsilon_{H,Z}) \le g(\mathbb{E}\,\epsilon_{H,Z}) \le g(\epsilon)$, whence
\[
  \Adv^{\extp}_{\calK,\calD,\calR}(S,\mathsf{D})
  \;\le\; \frac{1}{\sqrt2}\sqrt{\epsilon R}
  \;+\; \sqrt{\frac{\ln(2eD\epsilon)}{2K}} \;+\; \frac{9}{10\sqrt K} .
\]
For the closed form of \eqref{eq:main-plain}, note first that an $\epsilon$-unpredictable source exists only for
$\epsilon \ge 1/D$, since $\max_x p^{H,z}_x \ge 1/D$ pointwise and hence
$\mathbb{E}[\epsilon_{H,Z}] \ge 1/D$; in particular $D = 1$ forces $\epsilon = 1$. Put
$P := \epsilon R$, $L := \log_2 D$, and
$f(P,L) := \bigl(\tfrac{1}{\sqrt2}\sqrt P + \tfrac65\sqrt{(1+(\ln2)L)/K}\bigr)/\sqrt{P+L/K}$, the
ratio of the two bounds. If $D = 1$ then $L = 0$ and $P = \epsilon R = R \ge 1$, since $\epsilon = 1$,
so $f = \tfrac{1}{\sqrt2} + \tfrac65 (KP)^{-1/2} \le \tfrac{1}{\sqrt2} + \tfrac65 < 1.91$. If
$D \ge 2$ then $L \ge 1$, and Fact~\ref{fact:cs} applied to the pairs
$\bigl(\tfrac{1}{\sqrt2},\, \tfrac65\sqrt{(1+(\ln2)L)/L}\bigr)$ and $(\sqrt P, \sqrt{L/K})$ gives
\[
  f \;\le\; \sqrt{\tfrac12 + \tfrac{36}{25}\Bigl(\tfrac1L + \ln 2\Bigr)}
    \;\le\; \sqrt{\tfrac12 + \tfrac{36}{25}(1+\ln 2)} \;<\; 1.72 .
\]
So $f < 1.91$ whenever $D \ge 2$, for every $\epsilon \in (0,1]$, the step of Fact~\ref{fact:cs}
having used only $P \ge 0$; and $f < 1.91$ for $D = 1$, $\epsilon = 1$. In particular the constant
$c = 2$ of the corrected conjecture suffices throughout the satisfiable range. (For $D = 1$ the
restriction to $\epsilon = 1$ is necessary: for $D = R = 1$ and $\epsilon$ small the second
inequality fails.) This proves both inequalities of \eqref{eq:main-plain} and the inequality
\eqref{eq:main-sharp}, which together are Theorem~\ref{thm:main}.
\hfill$\square$

\section{Discussion}\label{sec:disc}

Five things are worth separating: what the public seed costs and what it leaves intact
(Remark~\ref{rem:cost}); where the shape of the bound comes from, and which part of the hidden-seed
intuition does not survive the change (Remark~\ref{rem:shape}); how far the two summands are
matched from below, and which half of that is proved rather than sketched
(Remark~\ref{rem:tight}); which of \eqref{eq:main-plain} and \eqref{eq:main-sharp} to prefer at
given $D$ and $\epsilon$ (Remark~\ref{rem:refine}); and, last, which three features of the model
the proof actually consumes.

\begin{remark}[What the public seed costs, and what it does not]\label{rem:cost}
Comparing with the hidden-seed bound of Fact~\ref{fact:hidden}, exactly one factor $1/K$ is lost, and it is lost from the first summand only. The reason is
visible in Lemma~\ref{lem:analytic}: the average over seeds moves from inside the statistical
distance to outside it. Inside, it mixes $K$ rows into one distribution supported on up to $K/\epsilon$
values, and the mixture is smooth; outside, each row is judged on its own $1/\epsilon$ values and
no amount of averaging over rows repairs a row that is individually far from uniform. Accordingly
the available entropy drops from $\log_2 K + \log_2(1/\epsilon)$ to $\log_2(1/\epsilon)$, and
$\tfrac{1}{\sqrt2}\sqrt{\epsilon R}$ is the classical leftover-hash shape with an entropy loss of
$2\log_2(1/\epsilon)$ bits. The selection term keeps its $1/K$ in full, by
Lemma~\ref{lem:unif}. So the correct summary is not that the seed becomes useless: it becomes
useless against entropy deficiency, and remains exactly as effective against adversarial selection.
\end{remark}

\begin{remark}[Where the shape comes from]\label{rem:shape}
The step that decides the theorem is again the choice of a deviation $\lambda$ \emph{independent of
$t$} in Lemma~\ref{lem:unif}. The union-bound cost of ranging over supports of size $t$ is
$\ln\binom{D}{t} \approx t \ln D$, while $F_T(H)$ is sub-Gaussian with variance proxy $1/(4Kt)$.
In the hidden-seed model that variance proxy came from $Q_{H,T}$ averaging $Kt$ oracle values; here
it comes from $F_T$ averaging $K$ row-statistics with weight $1/K$ each, a statistic of $t$ values
having bounded-differences constant $1/t$. The arithmetic is the same and the two factors of $t$
cancel, so the required deviation is $\sqrt{t \ln D/(Kt)} = \sqrt{\ln D/K}$ for every $t$ at once,
and the selection term enters additively. What does \emph{not} carry over is the accompanying
intuition about min-entropy: it is no longer true that a source with more entropy averages away
more bias, since each row now sees only $t$ values. The cancellation here is between the
union-bound cost and the $1/K$ from averaging over rows, which is a statement about the seed space
rather than about the source. Nor does the term enter uniformly in $\epsilon$, as
\eqref{eq:main-sharp} shows: what cancels exactly is $t$ against $t$, and what survives is the
per-size cost $\tfrac1t\ln\binom Dt$, which is $\ln D$ only for $t \ll D$ and vanishes at $t = D$.
The uniform-in-$\epsilon$ reading is an artefact of estimating $\tfrac1t\ln\binom Dt$ by
$1 + \ln D$.
\end{remark}

\begin{remark}[Matching lower bounds, sketch]\label{rem:tight}
The proof gives $O(\sqrt{\epsilon R}) + O(\sqrt{\log D/K})$. About the first summand two different
things can be said, and it is worth separating them. That it must not decay in $K$ at all is
\emph{proved}, not sketched: Proposition~\ref{prop:false} exhibits, at the single point
$\epsilon R = 1$, a source with advantage at least $\tfrac14$ against
$\tfrac{1}{\sqrt2}\sqrt{\epsilon R} = 0.708$, uniformly in $K$. That the shape $\sqrt{\epsilon R}$
and the constant are right throughout $\epsilon R \le 1$ is a sketch, and one point cannot
establish it; the witness is instead the fixed-support flat source ($T$ fixed, $|T| = t = 1/\epsilon$,
$z = \bot$), for which $\Delta_{H,z} = F_{T}(H)$ exactly, so that
$\mathbb{E}_{H}[\Delta_{H,z}] = \mathbb{E}\,\SD(\emp^{H}_{k,T},U_{\calR})$ approaches
$\sqrt{2/\pi}\cdot\tfrac12\sqrt{R/t} = 0.3989\sqrt{\epsilon R}$ for $R \le t$, each
$|\emp^{H}_{k,T}(r) - 1/R|$ being asymptotically half-normal. Granting that estimate, the constant
of Lemma~\ref{lem:mean} is tight within $0.5/0.3989 = 1.25$ and the $\tfrac{1}{\sqrt2}$ of
\eqref{eq:main-plain} within $1.77$, a sharper comparison than the factor
$0.708/0.25 = 2.83$ that Proposition~\ref{prop:false} gives, but a sketch where the proposition is
a proof.
For the second we only sketch. Fix $A \subseteq \calR$ with $|A| = R/2$, let $S$ inspect the whole
table, let $T$ consist of the $1/\epsilon$ inputs $x$ maximising $|\{k : H(k,x) \in A\}|$, and let
$S$ output $x \sample T$ with $z \leftarrow T$. Since
$F_T(H) \ge \bigl|\tfrac{1}{K}\sum_{k}(\emp^{H}_{k,T}(A) - \tfrac12)\bigr|$, extremal-value
considerations over the $D$ columns give a bias of order $\sqrt{\log(D\epsilon)/K}$, which is of
order $\sqrt{\log D/K}$ only when $1/\epsilon \le D^{1-\Omega(1)}$: selecting the $t$ most biased of
$D$ columns buys $\sqrt{\ln(D/t)/K}$, not $\sqrt{\ln D/K}$, and at $t = D$ there is no selection
freedom at all. That $\sqrt{\log(D\epsilon)/K}$ is the shape \eqref{eq:main-sharp} matches, and it is why we state that
bound alongside \eqref{eq:main-plain}; whether the constants can be brought together we do not
settle here.
\end{remark}

\begin{remark}[Which of the two bounds to use]\label{rem:refine}
Neither of \eqref{eq:main-plain} and \eqref{eq:main-sharp} implies the other, since the additive
$\tfrac{9}{10}K^{-1/2}$ of the second is the price of dropping $\ln D$ in favour of
$\ln(2eD\epsilon)$. Both selection terms are of the form $\gamma/\sqrt K$, so the comparison is
between the two constants and is independent of $K$ and $R$. At $\epsilon = 1$ the constants are
$\tfrac65\sqrt{1+\ln D}$ and $\sqrt{\ln(2eD)/2} + \tfrac9{10}$, and the second is smaller from
$D \ge 26$ onwards (at $D = 25$ they are $2.4648$ against $2.4672$, at $D = 26$ they are $2.4762$
against $2.4734$); asymptotically it is smaller by a factor $\tfrac65\sqrt2 \approx 1.70$. At
maximal entropy $\epsilon = 1/D$ the second constant is $\sqrt{\ln(2e)/2} + \tfrac9{10} < 1.83$
\emph{whatever $D$ is}, while the first grows without bound, and the crossover is already at
$D = 4$. (Retaining the sharp constant $\tfrac32\sqrt{\ln2/2} = 0.8831$ of
Lemma~\ref{lem:unif}(ii) in place of $\tfrac9{10}$ moves these two thresholds to $D \ge 23$ and
$1.81$.) The reason is the one recorded in Remark~\ref{rem:tight}: at $\epsilon = 1/D$ the only
admissible support is $T = \calD$, there is nothing to select, and a selection term growing in $D$
is then an artefact of the estimate $\binom{D}{t} \le e^{ta}$ rather than a real cost. Everything
in this remark and in part~(ii) of Lemma~\ref{lem:unif} applies verbatim to the hidden-seed proof,
where it likewise replaces $1+\ln D$ by $\ln(2eD\epsilon)$.
\end{remark}

\begin{remark}[What the proof does and does not use]
Three features of the model are used, and only these. (a) $\sd$ is sampled after $S$ terminates and
is therefore independent of $(X,Z,H)$: this is what makes the same $p^{H,z}$ appear in every row's
pushforward $\nu^{H,z}_{k}$, and it supplies the $1/K$ in the selection term. Secrecy of $\sd$ is
\emph{not} used anywhere, which is why the argument transfers; independence is used everywhere. (b)
The unpredictability hypothesis is used only through $\mathbb{E}_{(H,Z)}[\max_x p_x] \le \epsilon$,
and only via Lemma~\ref{lem:flat} and one application of Fact~\ref{fact:jensen}; no pointwise bound on
$\epsilon_{H,z}$ is needed. (c) The source's use of $z$ is never analysed: since
Lemma~\ref{lem:unif} is uniform over \emph{all} subsets of \emph{all} sizes, it covers every
measurable rule by which $S$ may pick a support from $H$ and signal it through $z$. Consequently the
theorem holds verbatim for every $S$ whose auxiliary output is an arbitrary function of $H$ and its
coins, including $z = H$. Lemma~\ref{lem:unif} also yields a high-probability refinement, but only
for the selection component: with probability $1 - O(e^{-u})$ over $H$ alone, the deviations
$W_{1}$ and $W_{2}$ are small for all supports, hence all sources, simultaneously. No analogous per-$H$ statement holds for
the deficiency component: the hypothesis bounds $\epsilon_{H,Z}$ only on average over $(H,Z)$, and
an $\epsilon$-unpredictable source may concentrate its predictable mass on an arbitrary fixed $H$-event of
measure $\epsilon$, so the full advantage bound is inherently an average over $H$.
\end{remark}
\cjpart{II}{What the Leftover Hash Lemma Gives}
\runninghead{WHAT THE LEFTOVER HASH LEMMA GIVES}

Part~I settles the conjecture. Part~II asks what settling it is worth. Set beside the bounds already
known for extraction from a random oracle, what does the conjectured bound buy? It buys seed length,
and by a wide margin. To certify a target $\delta$ the collision route needs about $\log_2 R$ bits
of seed beyond the $2\log_2(1/\delta)$ every route pays. A $2$-universal family needs $D/R$ seeds,
so at these lengths none exists. The unsalted random-oracle heuristic needs no seed, but a source
that reads the oracle defeats it outright. Bounding the adversary's queries instead costs
$\log_2 q$ bits. The conjecture asks for $\log_2\log_2 D + 2\log_2(1/\delta) + 2$, which is $77$
bits in the leading example against $190$, $1920$ and $113$. It is also the only one of the four
whose seed grows with none of these: the derived key's length, the adversary's budget, and the
secret's length beyond a doubly-logarithmic term.

Nothing further is resolved here. The shape of this part is: notation, then every bound stated with
no proof, then the comparison, then what is not claimed. A reader who wants only the results needs
Sections~\ref{sec:budgets} to~\ref{sec:compare}; the proofs and the rest are in the appendices, all
of which belong to this part. One thing here is not new: the floor the collision route runs into is
the leftover hash lemma's known seed-length limit, which Barak et al.\ attribute to Stinson
(Appendix~\ref{sec:prior}). What is new is the comparison.

\section{Query budgets and collision notation}\label{sec:budgets}

The game is the $\Gm\ \Extp$ of Section~\ref{sec:setting} throughout, and $\epsilon$-unpredictability
is Definition~\ref{def:games}'s, which by Lemma~\ref{lem:analytic} says
$\mathbb{E}_{(H,Z)}[\epsilon_{H,Z}] \le \epsilon$. Two abbreviations run from here on: the parameter
subscripts are dropped, so that $\Adv^{\extp}(S,\mathsf{D})$ means
$\Adv^{\extp}_{\calK,\calD,\calR}(S,\mathsf{D})$; and a table is written $h$ rather than $H$ where it
is held fixed. What this part adds to the setting is a query model and one collision statistic.

\medskip
\noindent\textbf{Query budgets.} Later sections also treat the variant in which each party gets
oracle access in place of the table. There $S^{H}$ makes at most $q_S$ queries before returning
$(x,z)$, the predictor at most $q_P$, and $\mathsf{D}^{H}(\sd,y_b,z)$ at most $q_D$. The seed is
still drawn after $S$ stops. Write $Q_S \subseteq \calK\times\calD$ for the set of cells the source
queries. Two features of a source govern what follows, and neither is $q_S$:
\begin{equation}\label{eq:rs-ss}
  r_S \;:=\; \sum_{k \in \calK} \Pr\bigl[Q_S \cap (\{k\}\times\calD) \neq \emptyset\bigr] ,
  \qquad
  s_S \;:=\; \sum_{k \in \calK} \Pr\bigl[(k,X) \in Q_S\bigr] ,
\end{equation}
the expected number of \emph{rows} the source touches and the expected number of seeds at which it
queries \emph{its own output}. Both are at most $\min(q_S,K)$, and both can be far smaller than
$q_S$.

Unqualified, $\epsilon$ means the strongest reading: the predictor is unbounded and holds the whole
table. Every bound in Part~II assumes it except Theorem~\ref{thm:badcell}, which needs only a predictor
as rich as the distinguisher. Appendix~\ref{sec:predictor} shows that hypothesis is forced.

\medskip
Throughout, condition on a table and an auxiliary output, $(H,Z) = (h,z)$, and write
$p_x := \Pr[X = x \mid h, z]$ for the source's conditional law and $\epsilon_{h,z} := \max_x p_x$.
The quantities
\[
  \CP \;:=\; \sum_{x \in \calD} p_x^2 ,
  \qquad
  c_k \;:=\; \Pr_{X,X'}\bigl[h(k,X) = h(k,X')\bigr] ,
  \qquad
  \bar c \;:=\; \frac1K \sum_{k \in \calK} c_k ,
\]
all depend on $(h,z)$, which the notation suppresses. Inside an expectation over $(H,Z)$ read them
as the matching random variables. Here $X, X'$ are independent draws from $p$. Note
$\CP \le \epsilon_{h,z}$, since $\sum_x p_x^2 \le (\max_x p_x)\sum_x p_x$.

\section{The bounds}\label{sec:bounds}

Every bound the comparison uses, stated without proof. The proofs are in
Appendix~\ref{sec:proofs}. Two bounds are quoted rather than proved anywhere here: the conjecture,
which Part~I proves, and the results of \cite{BDKPPSY11}. Whatever does not bear on
seed length is in the appendices.

\subsection*{The conjecture}

Conjecture~\ref{conj:stmt-pub}, displayed as \eqref{eq:conj} in Part~I, asserts a universal
constant $c$ such that
$\Adv^{\extp}(S,\mathsf{D}) \le c\bigl(\sqrt{\epsilon R} + \sqrt{\log_2 D/K}\bigr)$
for every $\epsilon$-unpredictable $S$ and every unbounded $\mathsf{D}$. Every numerical
comparison below uses \eqref{eq:conj} exactly as displayed, with $c = 2$, the constant Part~I
establishes in Corollary~\ref{cor:pub}. No other constant from Part~I is imported and none is
needed.

\subsection*{Also stated, not proved here: the bounds of Barak et al.}

This note measures against several bounds of~\cite{BDKPPSY11}, so they are listed here with the
rest. All are quoted, not proved. All live in that paper's model rather than the game of
Section~\ref{sec:setting}: there $\mathrm{Ext}$ is a fixed public function, the source is a
distribution rather than an algorithm, and the random oracle appears in no definition, lemma or
theorem. Remark~\ref{rem:roquery} sets the two models side by side.

First a clash of notation. In~\cite{BDKPPSY11}, $v$ is the number of extracted bits, $m$ the
source's min-entropy, $L := m - v$ the \emph{entropy loss}, and $\varepsilon$ the security of the
\emph{application} being keyed; the source's length is written $|X|$. Part~II writes $u$ for
that length and $\mu$ for their $m$, since $m$ is needed below for $\log_2 R$; the renaming is
ours and not theirs. That last is not this paper's
$\epsilon$, which is the source's unpredictability and is $2^{-\mu}$ in their units. Otherwise
$u = \log_2 D$ and $v = \log_2 R$.

\begin{enumerate}
\item \textbf{Seed length} (their introduction, attributed to Stinson~\cite{Sti94}): a family
(almost) universal enough for the leftover hash lemma needs
$n \ge \min\bigl(u-v,\ v + 2\log(1/\varepsilon)\bigr) - O(1)$, ``and must grow with the number of
extracted bits''. Theorem~\ref{thm:limit} and Proposition~\ref{prop:nouniv} reproduce this, one
branch each; Section~\ref{sec:prior} makes the identification.

\item \textbf{Entropy loss} (their introduction): the standard lemma needs
$L \ge 2\log(1/\varepsilon)$.

\item \textbf{Their generalised leftover hash lemma} (Theorem~3.1, applied to particular primitives
in their Theorems~3.3--3.7): the loss falls to $L = \log(1/\varepsilon)$ for a wide range of
key-derivation applications, the security degrading from $\varepsilon$ to at most
$\varepsilon + \sqrt{\varepsilon\,2^{-L}}$.

\item \textbf{Expand-then-extract} (their Section~4): stretching a short public seed with a
pseudorandom generator and extracting with the result does \emph{not} work in general. They give a
DDH-based generator and a perfectly universal hash whose output is distinguishable with probability
close to $1$ given the short seed. It is sound when the number of extracted bits is logarithmic in
the generator's security, and sound in minicrypt, so any counterexample needs a public-key
assumption.

\item \textbf{The unsalted random-oracle heuristic} (their Section~3.4, prose rather than theorem):
$\varepsilon_{\mathrm{RO}} \approx T/2^{\mu} = \varepsilon\,2^{-L}$, for an adversary of running
time $T \approx \varepsilon 2^{v}$.
\end{enumerate}

Items 2, 3 and 5 are about \emph{entropy loss}; items 1 and 4 about \emph{seed length}. The two
axes are independent, and this part is on the second. Their entropy-loss improvement, item 3,
leaves item 1 where it was. Their own attempt on item 1 is item 4, and it fails in general. So the
seed-length limit is the more stubborn of the two, and it is the one Theorem~\ref{thm:limit} is
about.

\subsection*{What the collision route returns}

\begin{proposition}[The bound the lemma returns]\label{prop:lhl}
For every table $h$ and every view $z$ of positive probability,
\[
  \max_{\mathsf{D}} \; \bigl(2\Pr[\mathsf{D} \text{ wins} \mid h,z] - 1\bigr)
  \;\le\; \tfrac12\sqrt{R\,\bar c - 1} \, ,
\]
and consequently, for every unbounded $\mathsf{D}$,
$\Adv^{\extp}(S,\mathsf{D}) \le \mathbb{E}_{(H,Z)}\bigl[\tfrac12\sqrt{R\,\bar c - 1}\bigr]$.
\end{proposition}

Writing $c_k = \CP + \beta_k$ with
$\beta_k := \sum_{x \neq x'} p_x p_{x'} \mathbf{1}[h(k,x) = h(k,x')]$ and
$\bar\beta := \frac1K\sum_k \beta_k$ splits the quantity under the root as
\begin{equation}\label{eq:split}
  R \bar c - 1 \;=\; R\cdot\CP \;+\; \gamma ,
  \qquad \gamma \;:=\; R\bar\beta - 1 ,
\end{equation}
The first summand carries the source's entropy deficiency. The second carries the family's failure
to be universal on the source's support. A lower bound on $R\bar c - 1$ bounds only the sum, so
which summand it speaks about needs settling; Lemma~\ref{lem:attrib} settles it. Against a
universal table the second summand vanishes.

\begin{proposition}[Universal families]\label{prop:univ}
Fix a table $h$ and suppose the family $\{h(k,\cdot)\}_{k \in \calK}$ is $2$-universal, that is
$\Pr_{k}[h(k,x) = h(k,x')] \le 1/R$ for all $x \neq x'$. Then $\gamma \le -\CP \le 0$ for every
source law $p$, and the conditional bound of Proposition~\ref{prop:lhl} is at most
$\tfrac12\sqrt{R\,\epsilon_{h,z}}$. Consequently, if \emph{every} table in the support of $H$ has
this property, then $\Adv^{\extp}(S,\mathsf{D}) \le \tfrac12\sqrt{\epsilon R}$ for every
$\epsilon$-unpredictable $S$ and every unbounded $\mathsf{D}$.
\end{proposition}

The hypothesis is a property of a single table, whereas $H$ in Section~\ref{sec:setting} is uniform
on all of $\Fun(\calK\times\calD,\calR)$; the second sentence of Proposition~\ref{prop:univ} is
therefore about a different and easier model than the one at issue. It calibrates rather than
applies. On the first summand of \eqref{eq:conj} the leftover hash lemma gives
$\tfrac12\sqrt{\epsilon R}$, which is that summand with $c = \tfrac12$. Nothing is lost there, so
everything separating the two routes lies in $\gamma$.


\subsection*{Why that calibration does not apply, and what stands in its place}

Proposition~\ref{prop:univ} does not apply here. Universality is not just absent from a typical
table. At the seed lengths at issue no table has it, provided the output is shorter than the
input.

\begin{proposition}[No universal family at small $K$, when $R < D$]\label{prop:nouniv}
If $\{h(k,\cdot)\}_{k\in\calK}$ is $2$-universal then $D \le 2RK$; and, by Stinson's bound for
almost-universal families quoted in the proof, $D \le RK$, that is $K \ge D/R$. So no
$2$-universal family exists once $K < D/R$, which at $D = 2^{256}$ and $R = 2^{64}$ holds for every
$K < 2^{192}$. The proviso $R < D$ matters: for $R \ge D$ a single injective row is already
$2$-universal, with $K = 1$, so no seed-length lower bound holds without a hypothesis relating $R$
to $D$.
\end{proposition}

So $\gamma$ must be taken over every support the source may choose after reading $H$, and the next
theorem says how large those make $R\bar c - 1$. The source sees the whole table before $\sd$
exists, so any support it can name from $H$ is allowed.

\begin{theorem}[The collision route cannot beat $\sqrt{W/K}$]\label{thm:limit}
Let $\calK, \calD, \calR$ be nonempty and finite and let $\epsilon \in [1/D, 1]$. Put
\[
  t \;:=\; \lceil 1/\epsilon \rceil , \qquad
  \sigma \;:=\; \lceil Rt/D \rceil , \qquad
  W \;:=\; \lfloor R/\sigma \rfloor .
\]
Then $1 \le t \le D$, $1 \le \sigma \le R$ and $1 \le W \le R$. For \emph{every} table
$h \in \Fun(\calK\times\calD,\calR)$ there is a source $S$, uniform on a set of $t$ inputs, with
$z = \bot$ and $\epsilon$-unpredictable, for which
\[
  R\,\bar c - 1 \;\ge\; \frac{W - 1}{K} ,
  \qquad\text{so that the bound of Proposition~\ref{prop:lhl} is at least}\qquad
  \frac12\sqrt{\frac{W-1}{K}} \, .
\]
No randomness is used, so this holds for every table and not merely for most. When $R$ and $D/t$
are powers of two, as in the ranges of interest, $W = \min\{R,\ D/t\}$ exactly. In general
$W > \tfrac12\min\{R,\ D/t\} - 1$.
\end{theorem}

Theorem~\ref{thm:limit} bounds the sum $R\bar c - 1 = R\cdot\CP + \gamma$ from below. It says
nothing about $\gamma$ unless the floor exceeds what the first summand gives on its own, which for
this source is $R\cdot\CP = R/t$.

\begin{lemma}[Attribution]\label{lem:attrib}
For the source of Theorem~\ref{thm:limit}, $\CP = 1/t$ exactly, so
$\gamma \ge (W-1)/K - R/t$. In particular, if
\begin{equation}\label{eq:side}
  \frac{W-1}{K} \;>\; \frac{R}{t} ,
\end{equation}
then the floor of Theorem~\ref{thm:limit} is not explained by the entropy deficiency, and
$\gamma > 0$: the family is genuinely non-universal on that support, by a margin of at least
$(W-1)/K - R/t$.
\end{lemma}

Condition \eqref{eq:side} is not automatic, and the comparison below is stated only where it holds.
Take $R = D$ and $\epsilon = 1/2$, so $t = 2$, $\sigma = 2$, $W = \lfloor R/2\rfloor$, and $R/t = R/2$
exceeds $(W-1)/K$ for every $K \ge 1$. There the theorem is consistent with $\gamma \le 0$, that is
with a universal family, and says nothing about the second summand.


\subsection*{Bounding the queries}

Give the parties oracle access in place of the table. The leftover hash lemma cannot see the
distinguisher's budget at all. It bounds a statistical distance, and that is by definition the
supremum over every distinguisher, so pricing $q_D$ means leaving the lemma. Leaving it gives the
bound below, which is the one the seed comparison uses at a bounded budget. What becomes of the
collision route itself under query bounds is Appendix~\ref{sec:rowbreadth}; it does not change the
seed comparison.

\begin{theorem}[The bad-cell route]\label{thm:badcell}
For every $\epsilon$-unpredictable source and every distinguisher making at most $q_D$ queries,
\[
  \Adv^{\extp}(S,\mathsf{D}) \;\le\; \frac{s_S}{K} \;+\; q_D\,\epsilon
  \;\le\; \frac{\min(q_S,K)}{K} \;+\; q_D\,\epsilon .
\]
\end{theorem}

The first term of Theorem~\ref{thm:badcell} is not slack, and what it measures is the salt. The
next statement needs no probability and no queries.

\begin{proposition}[What the salt buys, and what its absence costs]\label{prop:unsalted}
For every table $h$ and every $K \ge 1$ there is a source, uniform on a set of at least $D/R$
inputs and so $\epsilon$-unpredictable with $\epsilon \le R/D$, with $r_S = s_S = 1$, and a
distinguisher making \emph{no} oracle query, such that
\[
  \Adv^{\extp}(S,\mathsf{D}) \;=\; \frac{1}{K}\Bigl(1 - \frac1R\Bigr) .
\]
In particular at $K = 1$ the advantage is $1 - 1/R$: unsalted extraction from a source that reads
the oracle fails outright, and it fails at min-entropy $\log_2 D - \log_2 R$, which is far above the
minimum any of the settings in Section~\ref{sec:compare} asks for. Since $s_S = 1$ and $q_D = 0$
here, Theorem~\ref{thm:badcell} returns $1/K$, so its first term is tight to within $1 - 1/R$.

More generally, take $1 \le j \le K$ and the source of Proposition~\ref{prop:qtight}, which pins $j$
rows at a uniform input and leaks their values. It has $s_S = j$ and is $\epsilon$-unpredictable
with $\epsilon \le R^{j}/D$, and against it a distinguisher making no oracle query achieves
\[
  \Adv^{\extp}(S,\mathsf{D}) \;=\; \frac{j}{K}\Bigl(1 - \frac1R\Bigr) .
\]
So the first term of Theorem~\ref{thm:badcell} is attained at every value of $s_S$, not only at
$s_S = 1$; and each pinned row costs $\log_2 R$ bits of min-entropy, exactly and not merely at
most, so a source required to retain min-entropy $\mu$ can afford only
$j \le (\log_2 D - \mu)/\log_2 R$ of them.
\end{proposition}

\section{The comparison}\label{sec:compare}

\subsection*{The bounds side by side}

Every bound of Section~\ref{sec:bounds} in one place. The source is unbounded throughout except
where a column says otherwise, and $\epsilon$ is $\epsilon_\infty$ except in the last row.

\begin{center}\small
\begin{tabular}{@{}lllll@{}}
\toprule
 & source & predictor & disting. & the bound \\
\midrule
\eqref{eq:conj}, conjectured        & unbounded    & unbounded      & unbounded    & $c\bigl(\sqrt{\epsilon R} + \sqrt{\log_2 D/K}\bigr)$ \\
Prop.~\ref{prop:lhl}, the route     & unbounded    & unbounded      & unbounded    & $\tfrac12\sqrt{\epsilon R + \gamma}$ \\
Prop.~\ref{prop:univ}$^{\ast}$      & unbounded    & unbounded      & unbounded    & $\tfrac12\sqrt{\epsilon R}$ \\
Thm.~\ref{thm:limit}, a \emph{floor}& unbounded    & unbounded      & unbounded    & $\ \ge\;\tfrac12\sqrt{(W-1)/K}$ \\
Thm.~\ref{thm:qlhl}                 & $r_S$ rows   & unbounded      & unbounded    & $\tfrac12\sqrt{\epsilon R + 2R r_S/K}$ \\
Thm.~\ref{thm:badcell}              & $s_S$ cells  & $q_P \ge q_D$  & $q_D$ quer.  & $s_S/K + q_D\,\epsilon$ \\
\bottomrule
\end{tabular}
\end{center}

\noindent In the second row $\gamma$ stands for $\mathbb{E}_{(H,Z)}[\gamma]$; the pointwise
statement is $\tfrac12\sqrt{R\cdot\CP + \gamma}$, and the displayed form is what Jensen's
inequality returns from it.

\smallskip
\noindent$^{\ast}$ On a \emph{fixed} $2$-universal table, which by Proposition~\ref{prop:nouniv}
needs $K \ge D/R$; the uniform $H$ of Section~\ref{sec:setting} is not one. Theorem~\ref{thm:qlhl}
and Corollary~\ref{cor:which} are stated and proved in Appendix~\ref{sec:rowbreadth}; they refine
the collision route but do not move the seed comparison, which is why they sit there.

\smallskip
Three readings of the table. The first two rows differ only in $\gamma$ against $\log_2 D/K$, and
Theorem~\ref{thm:limit} is what forbids $\gamma$ from being small enough to close that. The fifth
row is the fourth's matching ceiling for sources of bounded row breadth. And the last row is the
only one in which the distinguisher's budget appears at all, for the reason
Remark~\ref{rem:qd} gives; by Corollary~\ref{cor:which}, and at any $R$ these tables use, it is the
smallest of them whenever $q_D \le \tfrac12\sqrt{R/\epsilon}$.

Side by side, and this is the whole comparison in one line:
\begin{equation}\label{eq:sidebyside}
  \underbrace{c\Bigl(\sqrt{\epsilon R} \;+\; \sqrt{\tfrac{\log_2 D}{K}}\Bigr)}_{\text{conjectured, and proved in Part~I}}
  \qquad\text{against}\qquad
  \underbrace{\tfrac12\sqrt{\;\epsilon R \;+\; \tfrac{W-1}{K}\;}}_{\text{the least the collision route can return}} .
\end{equation}
Since $\sqrt{a+b} \ge (\sqrt a + \sqrt b)/\sqrt2$ by the arithmetic--geometric mean inequality, the
right-hand side is at least $\tfrac{1}{2\sqrt2}\bigl(\sqrt{\epsilon R} + \sqrt{(W-1)/K}\bigr)$, so
the two are the same two ingredients up to constants and differ in exactly one place: what stands
under the second square root. The conjecture pays $\log_2 D$, a number of bits; the collision route
pays $W$, the size of a set.


\subsection*{In entropy-loss units}

The three bounds of Barak et al.\ that share an axis, in that paper's own notation, with
$L = \mu - v$ the entropy loss and $\varepsilon$ the keyed application's security:

\begin{center}\small
\begin{tabular}{@{}llll@{}}
\toprule
 & security & min-entropy for $O(\varepsilon)$ & family, and status \\
\midrule
standard LHL                        & $2^{-L/2}$                    & $\mu = v + 2\log(1/\varepsilon)$ & universal; proved \\
\cite{BDKPPSY11}, Thm.~3.1          & $\sqrt{\varepsilon\,2^{-L}}$  & $\mu = v + \log(1/\varepsilon)$  & universal; proved \\
\cite{BDKPPSY11}, \S3.4             & $\varepsilon\,2^{-L}$         & $\mu = v$                        & unsalted RO; heuristic \\
\bottomrule
\end{tabular}
\end{center}

This is a different axis from the table above, and the two should not be read across. That one
varies the \emph{seed}, holding the family's universality fixed; this one varies the \emph{entropy
loss}, holding the seed fixed. Nothing in this column improves the seed length, which is the
quantity Theorem~\ref{thm:limit} constrains and the quantity the conjecture is about; and the last
row is not salted at all, so it has no seed to improve.

\subsection*{In seed length}

\begin{corollary}[The comparison, in seed length]\label{cor:seed}
Fix a target error $\delta \in (0,1)$ and parameters for which \eqref{eq:side} holds and the entropy
deficiency is not itself binding, that is $c\sqrt{\epsilon R} \le \eta\,\delta$ for some
$\eta \in (0,1)$. Then \eqref{eq:conj} certifies $\delta$ once
$K \ge c^{2}\log_2 D/\bigl((1-\eta)\delta\bigr)^{2}$, whereas by
Theorem~\ref{thm:limit} the leftover-hash route cannot certify $\delta$ at all unless
$K \ge (W-1)/(4\delta^{2})$. In bits,
\[
  \begin{aligned}
    \log_2 K \;&=\;    2\log_2 c + \log_2\log_2 D + 2\log_2(1/\delta)
                       + 2\log_2\tfrac{1}{1-\eta}                     && \text{(conjectured)}, \\
    \log_2 K \;&\ge\;  \log_2(W-1) - 2 + 2\log_2(1/\delta)            && \text{(collision route)} .
  \end{aligned}
\]
Both carry $2\log_2(1/\delta)$; they differ in the first term, $\log_2 W$ against
$\log_2\log_2 D$.
\end{corollary}

The last term of the first line is what the entropy deficiency costs the conjectured salt, and it
is the one place where the two lines are not the same kind of quantity: the second is a necessary
condition on $K$, the first a sufficient one, and driving \eqref{eq:conj}'s \emph{second} summand
alone below $\delta$ would not certify $\delta$. Everywhere below $\eta \le 2^{-15}$, so the term
is under $10^{-4}$ bits and the tables drop it.

\begin{remark}[A worked instance]\label{rem:numbers}
Take $D = 2^{256}$, $R = 2^{64}$, a source of min-entropy $160$ (so $\epsilon = 2^{-160}$ and
$t = 2^{160}$), $\delta = 2^{-32}$ and $c = 2$. Then $\sigma = \lceil 2^{-32}\rceil = 1$ and
$W = R = 2^{64}$. Both hypotheses of Corollary~\ref{cor:seed} hold: the deficiency term is
$c\sqrt{\epsilon R} = 2^{-47} \ll \delta$; and at the seed length \eqref{eq:conj} requires,
$K = 2^{74}$, the floor is $(W-1)/K \approx 2^{-10}$ against $R/t = 2^{-96}$, so \eqref{eq:side}
holds with room to spare and Lemma~\ref{lem:attrib} gives $\gamma > 0$. The comparison is then
$\log_2 K = 74$ under \eqref{eq:conj} against $\log_2 K \ge 126$ for the collision route: a gap of
$52$ bits.

At min-entropy $128$ instead, $c\sqrt{\epsilon R} = 2^{-31}$ exceeds $\delta = 2^{-32}$, so
\eqref{eq:conj} certifies that $\delta$ at no seed length and the comparison is empty. The
hypotheses of Corollary~\ref{cor:seed} are not decorative.
\end{remark}

\subsection*{In key-derivation units}

The game is a key-derivation problem. The input $x$ is a secret with entropy but no uniformity: a
Diffie--Hellman element, a biometric reading, the output of a physical source. The seed $\sd$ is a
public salt, $H$ is a public hash, and $H(\sd,x)$ is the derived key, which must look uniform to an
adversary who knows $H$ in full and sees $\sd$. In that reading $\log_2 K$ is the \emph{salt length}
in bits, $m := \log_2 R$ the derived key length, $n := \log_2 D$ the length of the secret's
encoding, and $\log_2(1/\epsilon)$ its min-entropy.

Whenever the derived key is no longer than the source's entropy gap, that is $m \le n - \log_2(1/\epsilon)$,
one has $W = R$ and the two salt lengths of Corollary~\ref{cor:seed} read
\[
  \underbrace{\log_2 n + 2\log_2(1/\delta) + 2}_{\text{conjectured}}
  \qquad\text{against}\qquad
  \underbrace{m + 2\log_2(1/\delta) - 2}_{\text{collision route}} .
\]
They differ in the first term, and they differ in kind. The conjectured salt grows with the
\emph{logarithm of the secret's length}. The collision route's grows with \emph{the length of the
key being derived}. So deriving a $256$-bit key rather than a $128$-bit one from the same secret
costs no extra salt under the conjecture, and $128$ further bits under the collision route. With
$c = 2$, and the hypotheses of Corollary~\ref{cor:seed} checked in each row:

\begin{center}\small
\begin{tabular}{@{}lrrrrrrrr@{}}
\toprule
 & \multicolumn{4}{c}{parameters} & \multicolumn{4}{c}{salt needed (bits)} \\
\cmidrule(lr){2-5}\cmidrule(l){6-9}
 & $\log_2 D$ & min-ent. & key & $\delta$ & conj. & route & univ. & RO$^{\dagger}$ \\
\midrule
DH element, 2048-bit group & $2048$ & $256$ & $128$ & $2^{-32}$ & $77$  & $190$ & $1920$ & $0$ \\
the same, 256-bit key      & $2048$ & $512$ & $256$ & $2^{-64}$ & $141$ & $382$ & $1792$ & $0$ \\
a 512-bit source           & $512$  & $300$ & $128$ & $2^{-64}$ & $139$ & $254$ & $384$  & $0$ \\
\bottomrule
\end{tabular}
\end{center}

\noindent\emph{The last column assumes the source is independent of $H$. It is the only column that
does, and Proposition~\ref{prop:unsalted} says what it becomes otherwise.}

\smallskip
\noindent$^{\dagger}$ The unsalted random-oracle heuristic of \cite{BDKPPSY11}, Remark~\ref{rem:roquery}:
$\epsilon_{\mathrm{RO}} \approx T/2^{\mu}$, so it needs no salt. That zero is the most informative
entry in the row, and it is bought rather than free. Two hypotheses buy it. The first is a
distinguisher of running time $T \le \delta\,2^{\mu}$, which is $T \le 2^{224}$, $2^{448}$ and
$2^{236}$ in the three rows. Those ceilings are enormous, which is why the heuristic is attractive
in practice. But they are ceilings, and all of the game of Section~\ref{sec:setting} lives above them, where every
party reads the whole table and $T/2^{\mu}$ says nothing.

The second hypothesis is that the source is independent of $H$. That zero does not survive dropping
it, not approximately and not by a little. By Proposition~\ref{prop:unsalted}, at $K = 1$ a source
that reads $H$ and is uniform on the largest fibre of the single row, of min-entropy
$\log_2 D - \log_2 R$, is distinguished from uniform with advantage $1 - 2^{-m}$ by a distinguisher
making \emph{no query at all}. That source is allowed in every row of the table: its min-entropy is
$1920$, $1792$ and $384$ against the $256$, $512$ and $300$ the rows ask for. Its entropy loss is
$L = \mu - v = 1792$, $1536$ and $256$, so the heuristic predicts
$\epsilon_{\mathrm{RO}} \le 2^{-1792}$, $2^{-1536}$ and $2^{-256}$ where the truth is
$1 - 2^{-128}$, $1 - 2^{-256}$ and $1 - 2^{-128}$. Off its hypothesis the estimate is not just
inapplicable. It is wrong by as much as an estimate of a probability can be.

So the first and last salt columns are the same primitive, a random oracle, priced under two
different assumptions about the source: nothing if it cannot look at $H$, and $77$ bits if it can.
The whole subject of this paper is the gap between those two entries.

The \emph{univ.} column is universal hashing, at $n - m$ bits from $K \ge D/R$
(Proposition~\ref{prop:nouniv}). It makes the point of Appendix~\ref{sec:vs} concrete: universal
hashing does not have a shorter salt than the conjecture, it has a far longer one. It is also read
differently from the first two. Those give the salt needed to reach the stated $\delta$; this gives
the salt needed for such a family to \emph{exist}. At that salt the error is
$\tfrac12\sqrt{\epsilon R}$, which is $2^{-65}$, $2^{-129}$ and $2^{-87}$ in the three rows, already
inside $\delta$ in each.

The three columns scale in three different ways, and that is the whole comparison in one place:
\[
  \underbrace{\log_2 n + 2\log_2(1/\delta) + 2}_{\text{conjectured}}, \qquad
  \underbrace{m + 2\log_2(1/\delta) - 2}_{\text{collision route}}, \qquad
  \underbrace{n - m}_{\text{universal}} .
\]
The conjectured salt grows with the logarithm of the secret's length. The collision route's grows
with the length of the derived key. Universal hashing's grows with the secret's length outright, and
\emph{shrinks} as the derived key grows: $1920$ bits for a $128$-bit key against $1792$ for a
$256$-bit one, since a larger $R$ weakens the constraint $K \ge D/R$. It does not depend on $\delta$
at all, being a requirement for existence rather than for accuracy. Only the first is what one wants
of a salt: short, and growing only in the logarithm of anything.

Read the other way, at a fixed salt rather than a fixed error, the same gap says the route
certifies almost nothing at the salt the conjecture licenses. Substituting
$K = c^{2}\log_2 D/\delta^{2}$ into Theorem~\ref{thm:limit}'s floor gives
\[
  \frac12\sqrt{\frac{W-1}{K}} \;\approx\; \frac{\delta}{2c}\sqrt{\frac{W}{\log_2 D}} ,
\]
so the error the route can certify exceeds $\delta$ by a factor of about
$\sqrt{W/\log_2 D}\,/\,(2c)$, roughly $(m - \log_2 n)/2$ bits of security, just under half the
derived key's length. In the first row that factor is $2^{56.5}$, so at a $77$-bit salt the route certifies
$2^{24.5}$, which is larger than $1$ and so says nothing. In the third row it certifies $2^{-6.5}$
at a $139$-bit salt: six bits of security where sixty-four were wanted.

Neither reading says that salting with $77$ bits is unsafe. Both say that the collision route cannot
be what establishes it, and Remark~\ref{rem:cause} is the reminder that the true advantage against
the witness source is far below either bound.

\subsection*{At a bounded query budget}

Suppose the source and the distinguisher each make at most $2^{n}$ oracle queries. Then
Theorem~\ref{thm:badcell} certifies $\delta$ once $s_S/K \le \delta/2$ and $q_D\epsilon \le \delta/2$,
and in the worst case over $2^{n}$-query sources $s_S \le 2^{n}$, so
\begin{equation}\label{eq:qsalt}
  \log_2 K \;\ge\; n + \log_2(1/\delta) + 1 ,
  \qquad\text{feasible only if}\qquad
  n \;\le\; \mu - \log_2(1/\delta) - 1 ,
\end{equation}
the second condition being about the source's min-entropy $\mu$ and not about the salt: no seed
length repairs a budget that already exceeds it. The three settings of the table above, alongside
the conjectured salt, which does not depend on $n$:

\begin{center}\small
\begin{tabular}{@{}lrrrrr@{}}
\toprule
 & \multicolumn{4}{c}{salt (bits) at $q_S = q_D = 2^{n}$} & conjectured \\
\cmidrule(lr){2-5}\cmidrule(l){6-6}
 & $n = 40$ & $n = 64$ & $n = 80$ & $n = 128$ & any $n$ \\
\midrule
DH element, 2048-bit group & $73$  & $97$  & $113$ & $161$ & $77$  \\
the same, 256-bit key      & $105$ & $129$ & $145$ & $193$ & $141$ \\
a 512-bit source           & $105$ & $129$ & $145$ & $193$ & $139$ \\
\bottomrule
\end{tabular}
\end{center}

Two things to read off. The first is that the last column is \emph{flat}: the conjectured salt is
$\log_2\log_2 D + 2\log_2(1/\delta) + 2$ whatever the adversary's budget, and it holds against an
unbounded one, whereas \eqref{eq:qsalt} grows one bit for every doubling of the budget. They cross
at $n = 44$, $76$ and $74$ in the three rows. So at $2^{80}$, the budget one would actually quote,
bounding the queries buys nothing: it asks for $113$ bits where the conjecture asks for $77$, and
for $145$ against $141$ and $139$. The query bound helps only against adversaries weaker than about
$2^{44}$. The reason is in the two exponents: the bad-cell route pays $\log_2 q$ where the
conjecture pays $\log_2\log_2 D$, and only halves the $\delta$ term in exchange.

The second is that \eqref{eq:qsalt} is an upper bound on what is needed, and a generous one, because
$s_S \le 2^{n}$ is not achievable at these min-entropies. By Proposition~\ref{prop:unsalted} the
advantage $s_S/K$ is attained by pinning $s_S$ rows at the source's own output. But each pinned row
costs $\log_2 R$ bits of min-entropy, so the three rows admit only
$j \le (\log_2 D - \mu)/\log_2 R = 14$, $6$ and $1$ pinned rows. Against the best attack known the
salt actually forced is $\log_2 j + \log_2(1/\delta) = 36$, $67$ and $64$ bits, and
Proposition~\ref{prop:unsalted} on its own already forces $\log_2(1/\delta) = 32$, $64$ and $64$. So
for the query-bounded route the truth lies in $[36,113]$, $[67,145]$ and $[64,145]$, and the
conjectured $77$, $141$ and $139$ sit inside all three.

The slack has a source. A source may query $(k,x)$ at its own output for a great many $k$ and simply \emph{ignore} the
answers: that costs it no entropy and gains it no advantage, yet it inflates $s_S$ to the query
budget, and Theorem~\ref{thm:badcell} charges for it. A bound that charged only for queries whose
answers reach the distinguisher would close most of the gap; this paper does not prove one.


\section{What is not claimed}\label{sec:not}

Theorem~\ref{thm:limit} is one-sided. It says the collision route cannot return a bound better than
$\tfrac12\sqrt{(W-1)/K}$; it does \emph{not} say the route attains that. Every statement in this
note about the route's cost is therefore a necessary condition on $K$, never the cost itself, and
Corollary~\ref{cor:seed} is phrased that way deliberately. Theorem~\ref{thm:qlhl} supplies a
matching ceiling for sources of bounded row breadth, and Remark~\ref{rem:window} closes the window
to a factor of two on the witness itself; for sources that touch many rows the paragraph below still
stands.

The remaining window is wide. One can certify $\gamma$ for a random table by a bounded-differences
argument: $\bar\beta$ has bounded-differences constant $2/(Kt)$ in the $Kt$ table entries above a
support of size $t$, and a union bound over the $\binom{D}{t}$ supports of that size costs
$e^{t(1+\ln D)}$. For sources uniform on a support that yields $\gamma \le R\sqrt{2(1+\ln D)/K}$
with high probability, hence an upper bound of
$\tfrac12\sqrt{\epsilon R} + \tfrac12\sqrt R\,(2(1+\ln D)/K)^{1/4}$. Two gaps in that sketch are
worth naming. It covers sources uniform on a support only, and a general source needs a further
argument reducing to those. And the resulting ceiling exceeds Theorem~\ref{thm:limit}'s floor by a
factor of order $\sqrt{R/W}\,(K\ln D)^{1/4}$, about $2^{21}$ in the instance of
Remark~\ref{rem:numbers}. So for a source of unrestricted row breadth the truth lies somewhere in
that window, and locating it is not attempted.

Nor is the bound one would most like in the query-bounded setting of Section~\ref{sec:bounds}, namely
$s_S/K + \tfrac12\sqrt{\epsilon R}$: the pinned cells charged once each, and the leftover hash lemma
applied to what is left as though it were universal. That would match the true advantage against the
witness of Theorem~\ref{thm:limit} in both summands, and it is what one expects, since a row the
source touched away from its own output is a row it has learnt nothing exploitable about. Proving it
needs Lemma~\ref{lem:fresh} refined from ``row $k$ untouched'' to ``the single cell $(k,X)$
untouched'', and that refinement is not available by the argument given: the freshness of one cell
survives conditioning on the table only through the row-versus-rest factorisation the lemma uses,
and two runs of the source may query in each other's output columns. Theorem~\ref{thm:qlhl} and
Theorem~\ref{thm:badcell} therefore stand as two separate bounds rather than one combined one, and
Corollary~\ref{cor:which} takes their minimum instead of their better parts.

What the floor alone does establish, in the regime where \eqref{eq:side} holds and the deficiency
term is not binding, is that the route cannot reach \eqref{eq:conj} there: it needs at least
$\log_2(W-1) - 2\log_2 c - \log_2\log_2 D - 2$ further bits of seed, which is $52$ in the instance
of Remark~\ref{rem:numbers}. Outside that regime nothing is claimed. Remark~\ref{rem:cause} gives
parameters where the construction shows no loss at all: min-entropy $254$ over the same $D$ and
$R$.

Finally, the scope is one technique, not one lemma. Theorem~\ref{thm:limit} constrains the
particular functional $\tfrac12\sqrt{R\bar c - 1}$ of Proposition~\ref{prop:lhl}, arrived at by
bounding the collision probability of $(\sd, h(\sd,X))$ and passing to statistical distance through
\eqref{eq:cs-l1l2}. It says nothing about other uses of the leftover hash lemma, and the honest form
of the headline is that \eqref{eq:conj} is not a corollary of \emph{that route}. Nothing here is a
statement about \eqref{eq:conj} itself, which is true: Part~I proves it with $c = 2$, by
bounding the $L_1$ distance directly and never forming a collision probability at all.



\appendix

\section{Proofs}\label{sec:proofs}

Of the statements of Section~\ref{sec:bounds}, in the order they appear there. Two elementary
inequalities are recorded first, since several proofs use them.

\begin{fact}[Cauchy--Schwarz, two forms]\label{fact:csforms}
Let $q$ be a probability distribution on a finite set $\Omega$ and let $A \subseteq \Omega$ carry
all of its mass. Then
\begin{equation}\label{eq:cs-collide}
  \sum_{\omega} q_\omega^2 \;\ge\; \frac{1}{|A|} .
\end{equation}
If moreover $u$ is uniform on $\Omega$, then
\begin{equation}\label{eq:cs-l1l2}
  \SD(q,u) \;=\; \tfrac12 \|q - u\|_1 \;\le\; \tfrac12\sqrt{|\Omega|}\,\|q-u\|_2
  \;=\; \tfrac12\sqrt{|\Omega|\textstyle\sum_\omega q_\omega^2 - 1} \, .
\end{equation}
\end{fact}

\begin{proof}
For \eqref{eq:cs-collide}, $1 = (\sum_{\omega \in A} q_\omega)^2 \le |A| \sum_{\omega \in A}
q_\omega^2 \le |A| \sum_\omega q_\omega^2$. For \eqref{eq:cs-l1l2}, the inequality is
$\|v\|_1 \le \sqrt{|\Omega|}\,\|v\|_2$ applied to $v = q - u$, and the final equality is
$\|q-u\|_2^2 = \sum_\omega q_\omega^2 - 1/|\Omega|$, which follows by expanding the square and
using $\sum_\omega q_\omega = 1$.
\end{proof}
\begin{proof}[Proof of Proposition~\ref{prop:lhl}]
Because $\sd$ is drawn only after $S$ has terminated, it is uniform on $\calK$ and independent of
$X$ given $(h,z)$. The distinguisher's view therefore differs between the two challenge bits only
in the pair $(\sd, y_b)$, which is distributed as $Q := \mathrm{law}(\sd, h(\sd,X))$ under $b = 0$
and as the uniform distribution $U$ on $\calK \times \calR$ under $b = 1$. The optimal advantage of
an unbounded distinguisher between two distributions is their statistical distance, so the
left-hand side is $\SD(Q,U)$.

Now $|\calK \times \calR| = KR$ and
\[
  \sum_{k,r} Q(k,r)^2
  \;=\; \Pr\bigl[(\sd, h(\sd,X)) = (\sd', h(\sd',X'))\bigr]
  \;=\; \sum_{k} \frac{1}{K^2}\,c_k
  \;=\; \frac{\bar c}{K} ,
\]
the middle equality because $\sd,\sd'$ are independent and uniform, so both equal a given $k$ with
probability $1/K^2$, and given $\sd = \sd' = k$ the outputs collide with probability $c_k$. Feeding
this into \eqref{eq:cs-l1l2} gives
$\SD(Q,U) \le \tfrac12\sqrt{KR \cdot \bar c/K - 1} = \tfrac12\sqrt{R\bar c - 1}$, the radicand being
nonnegative by \eqref{eq:cs-collide} applied to $Q$ with $A = \calK\times\calR$. Averaging over
$(H,Z)$ gives the second claim.
\end{proof}
\begin{proof}[Proof of Proposition~\ref{prop:univ}]
Exchanging the order of summation,
\[
  \bar\beta \;=\; \sum_{x \neq x'} p_x p_{x'} \Pr_k[h(k,x) = h(k,x')] \;\le\; (1 - \CP)/R ,
\]
valid for every $p$ because the hypothesis is a statement about every pair $x \neq x'$; so
$\gamma \le 1 - \CP - 1 = -\CP$. By \eqref{eq:split} the conditional bound is at most
$\tfrac12\sqrt{R\cdot\CP - \CP} \le \tfrac12\sqrt{R\,\epsilon_{h,z}}$, using $\CP \le
\epsilon_{h,z}$. For the last claim, average over $(H,Z)$ and apply Jensen's inequality to the
concave nondecreasing $\sqrt{\cdot}$. This is legitimate although $p$ depends on $H$, since the
conditional bound holds pointwise in $(h,z)$. Combine it with
$\mathbb{E}[\epsilon_{H,Z}] \le \epsilon$.
\end{proof}
\begin{proof}[Proof of Proposition~\ref{prop:nouniv}]
Viewing each input $x$ as the column vector $(h(k,x))_{k \in \calK} \in \calR^{K}$, two distinct
columns agree in at most $K/R$ coordinates, so the $D$ columns form a code of length $K$ over an
alphabet of size $R$ with minimum distance at least $(1 - 1/R)K$. The Plotkin bound in its
equality case, ``$d = (1-1/q)n$ implies $|C| \le 2qn$'', gives $D \le 2RK$.

The factor of two belongs to that route rather than to the truth. Stinson's lower bound for
$\varepsilon$-almost-universal families~\cite{Sti94},
$K \ge D(R-1)/\bigl(\varepsilon R(D-R) + R^{2} - D\bigr)$, reads $K \ge D/R$ at
$\varepsilon = 1/R$, and that is the constant Section~\ref{sec:compare} uses. It is attained: over
$\mathbb{F}_q$ the family $h(k,(a,b)) = ak+b$ has $D = q^{2}$ and $R = K = q$. Reading the
elementary $D \le 2RK$ instead would shorten the \emph{univ.}\ column of
Section~\ref{sec:compare} by one bit and change nothing else.

For the proviso, if $R \ge D$ take
$K = 1$ and any injection $\calD \to \calR$: no two distinct inputs collide, so the family is
$2$-universal. (An earlier version of this material asserted $K \ge D^{\Omega(1)}$ unconditionally,
which that example refutes; and quoted the Plotkin bound in the sharper form $|C| \le qn$, which
holds for $q = 2$ but is not the general statement.)
\end{proof}
\begin{proof}[Proof of Theorem~\ref{thm:limit}]
\emph{The parameters are in range.} From $\epsilon \le 1$ we get $t \ge 1$, and from
$\epsilon \ge 1/D$ that $t = \lceil 1/\epsilon\rceil \le D$; hence $Rt/D \le R$ and, $R$ being an
integer, $\sigma = \lceil Rt/D\rceil \le R$, while $\sigma \ge 1$ because $Rt/D > 0$. From
$\sigma \le R$ we get $W = \lfloor R/\sigma\rfloor \ge 1$, and from $\sigma \ge 1$ that $W \le R$.

\emph{The support.} Fix any seed $k_1 \in \calK$ and let $\calS \subseteq \calR$ consist of
$\sigma$ symbols $r$ maximising $|\{x \in \calD : h(k_1,x) = r\}|$. The $D$ inputs are distributed
among the $R$ symbols, so the $\sigma$ most popular of them carry at least a $\sigma/R$ fraction of
the total, that is at least $D\sigma/R$ inputs; and $\sigma \ge Rt/D$ gives $D\sigma/R \ge t$. So
$\{x : h(k_1,x) \in \calS\}$ has at least $t$ elements; let $T$ be any $t$ of them and let $S(h)$
output $x \sample T$ together with $z \leftarrow \bot$. This is a counting fact about an arbitrary
map $\calD \to \calR$; no property of $h$ beyond being a function is used.

\emph{Admissibility.} $T$ is computed from $h$, which $S$ receives in full, and $\sd$ is drawn only
after $S$ terminates, so nothing in the construction depends on the seed and $S$ is a legitimate
source for the game. Its conditional law is uniform on $T$, so
$\epsilon_{h,\bot} = 1/t = 1/\lceil 1/\epsilon\rceil \le \epsilon$; since this holds for every
table, $\mathbb{E}[\epsilon_{H,Z}] \le \epsilon$ and $S$ is $\epsilon$-unpredictable in the sense of
Section~\ref{sec:setting}.

\emph{The two collision bounds.} For every $k$ the law of $h(k,X)$ is a distribution on $\calR$, so
\eqref{eq:cs-collide} with $A = \calR$ gives $c_k \ge 1/R$. For $k = k_1$ every $x \in T$ has
$h(k_1,x) \in \calS$, so the law of $h(k_1,X)$ is carried by $\calS$ and \eqref{eq:cs-collide} with
$A = \calS$ gives $c_{k_1} \ge 1/\sigma$.

\emph{Assembling.} Averaging the $K$ lower bounds and using $R/\sigma \ge \lfloor R/\sigma\rfloor
= W$,
\[
  R\,\bar c - 1
  \;\ge\; R\Bigl[\frac1K\cdot\frac1\sigma + \frac{K-1}{K}\cdot\frac1R\Bigr] - 1
  \;=\; \frac{R/\sigma}{K} + \frac{K-1}{K} - 1
  \;\ge\; \frac{W}{K} - \frac1K
  \;=\; \frac{W-1}{K} ,
\]
and Proposition~\ref{prop:lhl} applied to this source returns
$\tfrac12\sqrt{R\bar c - 1} \ge \tfrac12\sqrt{(W-1)/K}$.

\emph{The two forms of $W$.} If $R = 2^{m}$ and $D/t = 2^{\ell}$ then $Rt/D = 2^{m-\ell}$, so
$\sigma = \max\{1, 2^{m-\ell}\}$ and $W = \lfloor R/\sigma\rfloor = \min\{2^{m}, 2^{\ell}\}
= \min\{R, D/t\}$. In general $\sigma < Rt/D + 1$, so, writing $a := R$ and $b := D/t$,
\[
  \frac{R}{\sigma} \;>\; \frac{a}{a/b + 1} \;=\; \frac{ab}{a+b}
  \;=\; \frac{\min\{a,b\}\cdot\max\{a,b\}}{a+b}
  \;\ge\; \tfrac12\min\{a,b\} ,
\]
the last step because $a + b \le 2\max\{a,b\}$; hence
$W > R/\sigma - 1 > \tfrac12\min\{R, D/t\} - 1$.
\end{proof}
\begin{proof}[Proof of Lemma~\ref{lem:attrib}]
The source is uniform on $T$ with $|T| = t$, so $\CP = \sum_{x\in T} t^{-2} = 1/t$. Substituting
into \eqref{eq:split} and rearranging Theorem~\ref{thm:limit}'s conclusion gives
$\gamma = R\bar c - 1 - R/t \ge (W-1)/K - R/t$, and \eqref{eq:side} makes it positive.
\end{proof}
\begin{proof}[Proof of Theorem~\ref{thm:badcell}]
Write the advantage as $|\Pr[b' = 1 \mid b = 0] - \Pr[b' = 1 \mid b = 1]|$ and couple the two worlds
by sampling $H$ lazily, with the same coins for $S$ and $\mathsf{D}$ in both. They differ only in
the challenge: the cell $H(\sd,x)$ under $b = 0$, an independent uniform value under $b = 1$. Let
$\mathrm{BAD}$ be the event that the cell $(\sd,x)$ is queried, by the source or by the
distinguisher. Off $\mathrm{BAD}$ that cell is never read by anyone, so the two executions coincide
step for step: they can first diverge only at a query to $(\sd,x)$, whence $\mathrm{BAD}$ occurs in
one world exactly when it occurs in the other and the two views agree off it. The advantage is
therefore at most $\Pr[\mathrm{BAD}]$, which we may compute in the $b = 1$ world, and by a union
bound at most
$\Pr[(\sd,x) \in Q_S] + \Pr[\mathsf{D} \text{ queries } (\sd,x)]$.

For the first, $\sd$ is uniform on $\calK$ and independent of the source's run, so
$\Pr[(\sd,X) \in Q_S] = \frac1K\sum_k \Pr[(k,X) \in Q_S] = s_S/K$, and $s_S \le \min(q_S,K)$ because
the source can place at most $q_S$ queries and there are only $K$ rows.

For the second, build a predictor. Given $(H,z)$, let $P$ draw $\sd \sample \calK$ and
$y \sample \calR$, run $\mathsf{D}^{H}(\sd, y, z)$ answering its queries from $H$, pick
$i \sample \{1,\dots,q_D\}$ and return the input component of the $i$-th query. (At $q_D = 0$
the distinguisher queries nothing, the term is $0$, and no predictor is needed.) Under $b = 1$ the
challenge is a uniform value independent of everything else, so $P$'s simulation is exact; and if
$\mathsf{D}$ queries $(\sd,x)$ then some query has input component $x$, which $P$ selects with
probability at least $1/q_D$. So $\Pr[\mathsf{D}\text{ queries }(\sd,x)] \le q_D \cdot
\Adv^{\mathrm{pred}}(P) \le q_D\epsilon$.
\end{proof}
\begin{proof}[Proof of Proposition~\ref{prop:unsalted}]
Fix the seed $k_1 = 1$. Among the $R$ symbols, let $r^{*}$ be one maximising
$|h(1,\cdot)^{-1}(r^{*})|$, ties broken by index; the $D$ inputs are distributed among $R$ symbols,
so $T := h(1,\cdot)^{-1}(r^{*})$ has $|T| \ge D/R$ by pigeonhole. This is a counting fact about an
arbitrary map and uses no property of $h$. Let $S(h)$ read row $1$, compute $r^{*}$, and output
$x \sample T$ together with $z \leftarrow r^{*}$.

The source is admissible: $T$ is computed from $h$, which $S$ receives, and $\sd$ is drawn only
after $S$ terminates. Given $(h,z)$ its law is uniform on $T$, so
$\epsilon_{h,z} = 1/|T| \le R/D$ for every table, whence $\mathbb{E}[\epsilon_{H,Z}] \le R/D$.
Leaking $r^{*}$ costs nothing, since $z$ is a function of $h$ and the conditioning in
$\epsilon_{h,z}$ already fixes $h$. It queries only row $1$, and its output lies in that row's
queried cells, so $r_S = s_S = 1$.

Let $\mathsf{D}$ make no oracle query and return $1$ exactly when $\sd = 1$ and $y_b = z$. Under
$b = 0$ we have $y_0 = h(\sd,x)$, and when $\sd = 1$ every $x \in T$ gives $h(1,x) = r^{*} = z$, so
$\Pr[\mathsf{D} \Rightarrow 1 \mid b=0] = 1/K$. Under $b = 1$ the value $y_1$ is uniform on $\calR$
and independent of $z$, so $\Pr[\mathsf{D} \Rightarrow 1 \mid b=1] = 1/(KR)$. The advantage is the
difference, $(1 - 1/R)/K$.

For the general case, let $S$ draw $x \sample \calD$, query $(k,x)$ for $1 \le k \le j$ and output
$\bigl(x, z = (H(1,x),\dots,H(j,x))\bigr)$; Proposition~\ref{prop:qtight} gives
$\epsilon \le R^{j}/D$ and $s_S = j$, the latter because the source queries its own output in
exactly those $j$ rows. Let $\mathsf{D}$ make no query and return $1$ exactly when $\sd \le j$ and
$y_b = z_{\sd}$. Under $b = 0$, $y_0 = H(\sd,x) = z_{\sd}$ whenever $\sd \le j$, so the probability
is $j/K$; under $b = 1$ it is $j/(KR)$.

The entropy remark needs the bound $\epsilon \le R^{j}/D$ in both directions, so the upper bound
alone will not do. Proposition~\ref{prop:qtight}'s proof gives the equality
$\epsilon = \mathbb{E}_{H}\bigl[\#\{z : T_z \neq \emptyset\}\bigr]/D$; and a uniform $H$ leaves a
given one of the $R^{j}$ fibres empty with probability $(1-R^{-j})^{D} \le e^{-D/R^{j}}$, so all
$R^{j}$ are nonempty except with probability $R^{j}e^{-D/R^{j}}$, which is negligible whenever
$D/R^{j}$ comfortably exceeds $j\ln R$ --- it is $2^{256}$ against $1792\ln 2$ in the first row of
Section~\ref{sec:compare}, and larger in the other two. There $\epsilon = (1-o(1))R^{j}/D$, and
$\epsilon \le 2^{-\mu}$ rearranges to the stated ceiling on $j$.
\end{proof}

\section{The collision route under query bounds}\label{sec:rowbreadth}

Give the parties the budgets of Section~\ref{sec:budgets} instead of the table. The collision route
survives, in a form governed by row breadth rather than by query count. The source of
Theorem~\ref{thm:limit} reads an entire row, so $q_S = D$, yet $r_S = s_S = 1$.
\begin{theorem}[The collision route, query-bounded]\label{thm:qlhl}
For every source that is $\epsilon$-unpredictable and every unbounded distinguisher,
\[
  \Adv^{\extp}(S,\mathsf{D}) \;\le\; \frac12\sqrt{\;\epsilon R \;+\; \frac{2R\,r_S}{K}\;}
  \;\le\; \frac12\sqrt{\;\epsilon R \;+\; \frac{2R\min(q_S,K)}{K}\;} \, .
\]
\end{theorem}

Three readings of Theorem~\ref{thm:qlhl} are worth separating. It is a statement about \emph{row
breadth}, not query count: a source may spend $D$ queries inside one row and remain capped at
$\epsilon R + 2R/K$, while $K$ queries placed one per row exhaust the bound. It degrades to nothing
at $r_S = K$, correctly, since a source that pins every row leaves $\bar c = 1$. And at $r_S = 0$ it
returns $\tfrac12\sqrt{\epsilon R}$, the universal-family answer of
Proposition~\ref{prop:univ}: a source that never queries is independent of the table, and a random
table is universal in expectation.

\begin{proposition}[The dependence on $r_S$ is exact up to a factor of two]\label{prop:qtight}
Fix $j \le K$ and let the source draw $x \sample \calD$, query $(k,x)$ for $1 \le k \le j$, and output
$\bigl(x, z = (H(1,x),\dots,H(j,x))\bigr)$. It makes $j$ queries and has $r_S = s_S = j$; it is
$\epsilon$-unpredictable with $\epsilon \le R^{j}/D$; and for \emph{every} table
\[
  R\bar c - 1 \;\ge\; \frac{(R-1)\,j}{K} .
\]
\end{proposition}

So the term $2R\,r_S/K$ of Theorem~\ref{thm:qlhl} cannot be improved below $(R-1)r_S/K$: the
constant is off by $2R/(R-1)$, which is two in the ranges of interest, and the shape is right. At $j = K$ the proposition
returns $R - 1$, and the theorem is vacuous, as it must be.

\begin{remark}[No bound proved through Proposition~\ref{prop:lhl} can mention $q_D$]\label{rem:qd}
Proposition~\ref{prop:lhl} passes through $\SD(Q,U)$, and statistical distance \emph{is} the
supremum of the advantage over all distinguishers whatever. A bound obtained that way holds for the
unbounded distinguisher and can therefore never improve as $q_D$ falls. This is not a defect of the
present derivation but of the route: to let the distinguisher's budget count one must leave the
leftover hash lemma, and Theorem~\ref{thm:badcell} is what one gets by doing so.
\end{remark}
\begin{corollary}[Which route to use]\label{cor:which}
Under the hypotheses of both,
\[
  \Adv^{\extp}(S,\mathsf{D}) \;\le\;
  \min\Bigl\{\; \frac{s_S}{K} + q_D\epsilon , \;\;
    \frac12\sqrt{\epsilon R + \frac{2R\,r_S}{K}} \;\Bigr\} ,
\]
and if moreover $2s_S/K + 2\sqrt{\epsilon R} \le R$, the bad-cell route is the smaller whenever
\[
  q_D \;\le\; \tfrac12\sqrt{R/\epsilon} .
\]
\end{corollary}

The side hypothesis is not a real restriction: at $R \ge 2^{64}$, with $s_S \le K$ and
$\epsilon R \le 1$, its left-hand side is at most $4$. The threshold itself is not a close call
either. In the instance of Remark~\ref{rem:numbers} it is
$q_D \le 2^{111}$; for a $128$-bit derived key from a source of min-entropy $128$ it is
$q_D \le 2^{127}$. Below it, which is every parameter regime anyone would deploy, the leftover hash
lemma is the wrong instrument for this game. The gap is not marginal: on the witness of Theorem~\ref{thm:limit} in that instance, against a distinguisher making
$q_D = 2^{80}$ queries, Theorem~\ref{thm:badcell} returns $2^{-74} + 2^{-80} \approx 2^{-74}$ where
Theorem~\ref{thm:qlhl} returns $2^{-5.5}$.

\medskip\noindent\emph{Proofs for this appendix.} Theorem~\ref{thm:qlhl} needs one lemma of its own.

\begin{lemma}[An untouched row stays fresh]\label{lem:fresh}
Fix $k \in \calK$, let $A_k$ be the event that the source queries some cell of row $k$, and write
$p^{-}_{x} := \Pr[X = x \wedge \neg A_k \mid H, Z]$. Then
\[
  \mathbb{E}_{(H,Z)}\Bigl[\;\sum_{x \neq x'} p^{-}_{x}\,p^{-}_{x'}\,
    \mathbf{1}[H(k,x) = H(k,x')]\Bigr] \;\le\; \frac1R .
\]
\end{lemma}
\begin{proof}[Proof of Lemma~\ref{lem:fresh}]
Split the table as $H = (H_k, H_{-k})$, row $k$ against the rest; the two are independent and $H_k$
is uniform on $\Fun(\calD,\calR)$. Take $S$ deterministic given coins $\rho$, with $\rho$
independent of $H$.

Run $S$ on $(\rho, h)$. Its transcript up to its first query into row $k$ is a function of
$(\rho, h_{-k})$ alone; hence so is the indicator of $\neg A_k$, and on $\neg A_k$ so is the whole
transcript, and therefore $(X,Z)$. Write
\[
  a(x,z,h_{-k}) := \Pr_{\rho}[X = x \wedge Z = z \wedge \neg A_k \mid h_{-k}] ,
  \qquad
  A(z,h_{-k}) := \sum_{x} a(x,z,h_{-k}) ,
\]
so that $\Pr[H = h, Z = z, X = x, \neg A_k] = \Pr[h_{-k}]\Pr[h_k]\,a(x,z,h_{-k})$ and
\[
  \Pr[H = h, Z = z] \;=\; \Pr[h_{-k}]\Pr[h_k]\Pr_\rho[Z = z \mid h]
  \;\ge\; \Pr[h_{-k}]\Pr[h_k]\,A(z,h_{-k}) .
\]
Since $\Pr[H=h,Z=z]\,p^{-}_{x} = \Pr[H=h,Z=z,X=x,\neg A_k]$, the expectation in question is
\[
  \sum_{h,z} \frac{1}{\Pr[h,z]} \sum_{x \neq x'}
    \Pr[h,z,x,\neg A_k]\,\Pr[h,z,x',\neg A_k]\,\mathbf{1}[h(k,x) = h(k,x')] ,
\]
Substituting the two displays above, with the lower bound on $\Pr[h,z]$ in the denominator, which
is the direction an upper bound needs, makes it at most
\[
  \sum_{h_{-k}} \Pr[h_{-k}] \sum_{h_k} \Pr[h_k] \sum_{z} \frac{1}{A(z,h_{-k})}
    \sum_{x \neq x'} a(x,z,h_{-k})\,a(x',z,h_{-k})\,\mathbf{1}[h_k(x) = h_k(x')] .
\]
Neither $a$ nor $A$ depends on $h_k$, so the sum over $h_k$ may be performed first, and for
$x \neq x'$ it gives $\sum_{h_k}\Pr[h_k]\mathbf{1}[h_k(x) = h_k(x')] = 1/R$ exactly. What remains is
at most
\[
  \frac1R \sum_{h_{-k}}\Pr[h_{-k}] \sum_z \frac{A(z,h_{-k})^2}{A(z,h_{-k})}
  \;=\; \frac1R \sum_{h_{-k}}\Pr[h_{-k}]\,\Pr_\rho[\neg A_k \mid h_{-k}] \;\le\; \frac1R ,
\]
using $\sum_{x\neq x'} a\,a' \le (\sum_x a)^2 = A^2$ and $\sum_z A(z,h_{-k}) =
\Pr_\rho[\neg A_k \mid h_{-k}]$. Terms with $A(z,h_{-k}) = 0$ have vanishing numerator and are read
as $0$.
\end{proof}
\begin{proof}[Proof of Theorem~\ref{thm:qlhl}]
Fix $(h,z)$ and $k$, write $u_k := \Pr[A_k \mid h,z]$ and split $p = p^{+} + p^{-}$ into its parts on
$A_k$ and $\neg A_k$. Of the four terms of $c_k = \sum_{x,x'}p_x p_{x'}\mathbf{1}[h(k,x)=h(k,x')]$,
the three meeting $p^{+}$ contribute at most
$\sum_x p^{+}_x \sum_{x'} p_{x'} + \sum_{x'} p^{+}_{x'} = 2u_k$, the indicator being at most $1$.
Hence
\[
  c_k \;\le\; 2u_k \;+\; \sum_{x}(p^{-}_{x})^2 \;+\;
    \sum_{x \neq x'} p^{-}_{x}\,p^{-}_{x'}\,\mathbf{1}[h(k,x)=h(k,x')] ,
\]
and $p^{-} \le p$ pointwise makes the middle sum at most $\CP$. Average over $k$ and take the expectation over $(H,Z)$. For the first term,
$\mathbb{E}\bigl[\sum_k u_k\bigr] = \sum_k \Pr[A_k] = r_S$; for the second,
$\mathbb{E}[\CP] \le \mathbb{E}[\epsilon_{H,Z}] \le \epsilon$; for the third,
Lemma~\ref{lem:fresh} gives at most $1/R$ for each of the $K$ rows. So
\[
  \mathbb{E}[\,R\bar c - 1\,] \;\le\; \frac{2R\,r_S}{K} + \epsilon R + R\cdot\frac1R - 1
  \;=\; \epsilon R + \frac{2R\,r_S}{K} .
\]
By Proposition~\ref{prop:lhl} the advantage is at most
$\mathbb{E}\bigl[\tfrac12\sqrt{R\bar c - 1}\bigr]$, and $\sqrt{\cdot}$ is concave, so Jensen's
inequality moves the expectation inside the root. The radicand is nonnegative pointwise by
Fact~\ref{fact:csforms}, so this is legitimate.
\end{proof}
\begin{proof}[Proof of Proposition~\ref{prop:qtight}]
Given $(h,z)$ the conditional law is uniform on the joint fibre
$T_z := \{x' : h(k,x') = z_k \text{ for all } k \le j\}$, so $\epsilon_{h,z} = 1/|T_z|$, and since
$\Pr[Z = z] = |T_z|/D$ the average $\mathbb{E}[\epsilon_{H,Z}]$ is the number of nonempty fibres
divided by $D$, at most $R^{j}/D$. Every element of $T_z$ has the same image in each pinned row, so
$c_k = 1$ for $k \le j$; for the remaining rows $c_k \ge 1/R$ by \eqref{eq:cs-collide}. Hence
\[
  R\bar c - 1 \;\ge\; R\Bigl[\frac{j}{K}\cdot 1 + \frac{K-j}{K}\cdot\frac1R\Bigr] - 1
  \;=\; \frac{Rj}{K} - \frac{j}{K} \;=\; \frac{(R-1)j}{K} . \qedhere
\]
\end{proof}
\begin{proof}[Proof of Corollary~\ref{cor:which}]
Both bounds are proved above. Write $a := \epsilon R$, $\theta_s := s_S/K$ and
$\theta_r := r_S/K$, and note $\theta_s \le \theta_r$ because querying $(k,X)$ touches row $k$. At
$q_D \le \tfrac12\sqrt{R/\epsilon}$ one has $q_D\epsilon \le \tfrac12\sqrt a$, so it is enough that
$2\theta_s + \sqrt a \le \sqrt{a + 2R\theta_r}$. Squaring, that reads
$4\theta_s^{2} + 4\theta_s\sqrt a \le 2R\theta_r$, and $\theta_s \le \theta_r$ reduces it to
$4\theta_s + 4\sqrt a \le 2R$ when $\theta_s > 0$, which is the hypothesis; at $\theta_s = 0$ it is
immediate.

The two summands may \emph{not} be compared one at a time. It is true that
$s_S/K \le \tfrac12\sqrt{2Rr_S/K}$ always --- on squaring, $r_S \le RK/2$, which holds because
$r_S \le K$ and $R \ge 2$ --- but $\sqrt a + \sqrt b \ge \sqrt{a+b}$, so two termwise comparisons
do not add up to a comparison of the sums, and without a hypothesis the conclusion is false: at
$R = 2$, $r_S = s_S = K$, $\epsilon = 2^{-K}$ and $q_D = \tfrac12\sqrt{R/\epsilon}$ the bad-cell
bound is $1 + 2^{-(K+1)/2}$ against the collision route's $1 + 2^{-K-2}$ or so. An earlier version
of this material argued termwise.
\end{proof}

\begin{remark}[Theorem~\ref{thm:qlhl} closes the window of Section~\ref{sec:not}, for these sources]%
\label{rem:window}
Theorem~\ref{thm:limit} is a floor and Section~\ref{sec:not} names no matching ceiling; the one
sketched there exceeds the floor by a factor of about $2^{21}$ in the instance of
Remark~\ref{rem:numbers}. For a source of row breadth $r_S$ the two now meet. Take that instance,
$D = 2^{256}$, $R = 2^{64}$, min-entropy $160$, $K = 2^{74}$, where $W = R$. The witness of
Theorem~\ref{thm:limit} has $r_S = 1$, and
\[
  \frac{W-1}{K} \;\approx\; 2^{-10}
  \qquad\le\qquad R\bar c - 1 \qquad\le\qquad
  \epsilon R + \frac{2R}{K} \;=\; 2^{-96} + 2^{-9} \;\approx\; 2^{-9} :
\]
a window of a factor $2$ in the radicand, so the route returns between $2^{-6}$ and $2^{-5.5}$. The
hypotheses differ. Section~\ref{sec:not}'s sketch asks the source to be uniform on a support and
bounds no queries; Theorem~\ref{thm:qlhl} bounds $r_S$ and asks nothing about the shape. But the
witness satisfies both, and on it the query-bounded ceiling is sharper by a factor of about
$2^{20}$.
\end{remark}

\section{The predictor's budget}\label{sec:predictor}

The third party has a budget too, and it is not free to choose. Bounding the predictor
\emph{weakens} the hypothesis, since a predictor that cannot read the table is easier to defeat. So
the question is how weak the hypothesis may be made before the conclusions fail.
\begin{definition}[Unpredictability at a budget]\label{def:epsq}
For $q \in \{0,1,2,\dots\} \cup \{\infty\}$ let $\epsilon_q$ be the least $\epsilon$ such that
$\Pr[P^{H}(z) = X] \le \epsilon$ for every predictor $P$ making at most $q$ oracle queries. Then
$\epsilon_0 \le \epsilon_1 \le \cdots \le \epsilon_\infty$, and $\epsilon_\infty = \epsilon_{KD}$ is
the unbounded notion: a predictor allowed $KD$ queries may read the whole table, so
$\epsilon_\infty = \mathbb{E}_{(H,Z)}[\epsilon_{H,Z}]$.
\end{definition}
The two routes answer the question differently, and the difference is exactly the difference between
them. Note first that the reduction inside Theorem~\ref{thm:badcell} is query-preserving: the
predictor built there runs $\mathsf{D}$ and forwards each of its oracle queries, drawing $\sd$ and
$y$ itself, so it makes \emph{exactly} $q_D$ queries and never more. Its proof therefore establishes
the following without alteration.
\begin{theorem}[Three budgets]\label{thm:three}
Let $S$ make at most $q_S$ oracle queries, let $\mathsf{D}$ make at most $q_D$, and let the source
be $\epsilon$-unpredictable against predictors making at most $q_P$ queries. If $q_P \ge q_D$ then
\begin{equation}\label{eq:three-badcell}
  \Adv^{\extp}(S,\mathsf{D}) \;\le\; \frac{s_S}{K} \;+\; q_D\,\epsilon
  \qquad\text{with}\qquad s_S \le \min(q_S,K) ,
\end{equation}
and if $q_P = \infty$ then also
\begin{equation}\label{eq:three-lhl}
  \Adv^{\extp}(S,\mathsf{D}) \;\le\; \frac12\sqrt{\;\epsilon R + \frac{2R\,r_S}{K}\;}
  \qquad\text{with}\qquad r_S \le \min(q_S,K) .
\end{equation}
\end{theorem}

So all three budgets appear, and they appear in different places: $q_S$ through $r_S$ and $s_S$,
$q_D$ as a multiplier, and $q_P$ as a side condition on which bound may be invoked. Neither side
condition can be dropped, and one witness settles both.
\begin{proposition}[The predictor may not be cheaper than the distinguisher]\label{prop:qp}
Fix $R \le D/2$ and $1 \le m \le D/2$. Reserve the inputs $w_1,\dots,w_m$ from the upper half of
$\calD$ and identify $\calR$ with the lower half. Let the source query $(1,w_1),\dots,(1,w_m)$ and
output
\[
  x \;:=\; \Bigl(\sum_{j=1}^{m} H(1,w_j)\Bigr) \bmod R , \qquad z := \bot .
\]
Then $q_S = m$, $r_S = 1$ and $s_S = 0$; and
\begin{enumerate}
\item $\epsilon_q = 1/R$ for every $q < m$, while $\epsilon_q = 1$ for every $q \ge m$;
\item given $(H,Z)$ the source's law is a point mass, so $\CP = 1$, $c_k = 1$ for every $k$, and
      $R\bar c - 1 = R-1$;
\item an unbounded distinguisher, and equally one making $m+1$ queries, achieves advantage
      $1 - 1/R$.
\end{enumerate}
Consequently:
\begin{enumerate}
\item[(a)] In \eqref{eq:three-lhl} the hypothesis $q_P = \infty$ cannot be weakened to any finite
      budget. Taking $m = q_P + 1$ and $K \ge 2R$, the bound asserted with $\epsilon_{q_P}$ would be
      $\tfrac12\sqrt{1 + 2R/K} \le \tfrac12\sqrt2$, against a true advantage of $1 - 1/R$, which is
      larger for every $R \ge 4$.
\item[(b)] In \eqref{eq:three-badcell} the hypothesis $q_P \ge q_D$ cannot be weakened to
      $q_P \le q_D - 2$. Taking $m = q_D - 1$, the bound asserted with $\epsilon_{q_D-2}$ would be
      $q_D/R$, against a true advantage of $1 - 1/R$, which is larger whenever $R > q_D + 1$.
\end{enumerate}
\end{proposition}

The two halves of that proposition say the same thing about different routes. Reading them together:
\emph{the predictor must be granted at least what the distinguisher is granted, and granting it
exactly that is enough.} For the bad-cell route the budget is pinned to $q_D$ within a single query,
$q_P \ge q_D$ sufficing and $q_P \le q_D - 2$ failing. For the collision route the distinguisher is
unbounded by construction, by Remark~\ref{rem:qd}, so the predictor must be unbounded too. That is
why every bound of Part~II other than Theorem~\ref{thm:badcell} is stated against a predictor
holding the whole table. It is not a preference for the strongest hypothesis. It is forced.

\medskip\noindent\emph{Proofs for this appendix.}

\begin{proof}[Proof of Theorem~\ref{thm:three}]
For \eqref{eq:three-badcell}: $q_P \ge q_D$ gives $\epsilon \ge \epsilon_{q_P} \ge \epsilon_{q_D}$,
and the predictor of Theorem~\ref{thm:badcell} makes $q_D$ queries, so the bad event there is at
most $q_D\epsilon_{q_D} \le q_D\epsilon$. For \eqref{eq:three-lhl}, $q_P = \infty$ makes $\epsilon$
an upper bound for $\mathbb{E}[\epsilon_{H,Z}]$, which is the only use Theorem~\ref{thm:qlhl} makes
of the hypothesis.
\end{proof}
\begin{proof}[Proof of Proposition~\ref{prop:qp}]
The source queries only row $1$, so $r_S = 1$; its output lies in the lower half of $\calD$ and its
queries in the upper, so $(k,x) \notin Q_S$ for every $k$ and $s_S = 0$.

For (1), a predictor making $q < m$ queries leaves some $(1,w_i)$ unqueried, and conditioned on any
transcript that cell is uniform on $\calR$ and independent of the predictor's view; a sum modulo $R$
one of whose summands is uniform and independent is itself uniform, so $x$ is uniform on $\calR$
given that view and no guess does better than $1/R$. A predictor making $m$ queries reads the $m$
cells, which are fixed by the source's code, and computes $x$ outright.

For (2), $x$ is a function of $H$ alone, so the conditional law given $(H,Z)$ is a point mass:
$\CP = 1$ and $h(k,X) = h(k,X')$ always.

For (3), a distinguisher holding the table, or one spending $m$ queries on the reserved cells and
one more on $(\sd, x)$, computes $H(\sd,x)$ and returns $1$ exactly when $y_b$ equals it. Under
$b = 0$ it always returns $1$; under $b = 1$ it does so with probability $1/R$.

For (a), $m = q_P+1$ gives $\epsilon_{q_P} = 1/R$ by (1), so $\epsilon_{q_P}R = 1$; with $r_S = 1$
and $K \ge 2R$ the asserted radicand is at most $2$. For (b), $m = q_D - 1$ gives
$\epsilon_{q_D-2} = \epsilon_{m-1} = 1/R$ by (1), while a distinguisher making $m+1 = q_D$ queries
achieves $1 - 1/R$ by (3); the asserted bound is $s_S/K + q_D/R = q_D/R$.
\end{proof}

The witness also shows how little room there is between the two conditions. Its $x$ is a
deterministic function of $H$, so it is perfectly extractable-looking to any predictor poorer than
the source and worthless against any predictor as rich, with nothing in between: $\epsilon_q$ jumps
from $1/R$ to $1$ at $q = m$. The general form of that observation is worth recording.

\begin{remark}[Deterministic sources]\label{rem:det}
If $S$ is deterministic then $\epsilon_{q_S} = 1$, since a predictor with $q_S$ queries simply reruns
it. So for a deterministic source \eqref{eq:three-badcell} is vacuous once $q_D \ge q_S$, and
\eqref{eq:three-lhl} is vacuous outright. Extraction from such a source can only be secure against
distinguishers strictly poorer than the source itself, and the seed does not change that: a
distinguisher that can afford the source's own computation recovers $x$ and tests it with one
further query, whatever $K$ may be. Whatever entropy is being extracted has to sit in the source's
coins, not in the oracle it reads.
\end{remark}

\section{Further notes}\label{sec:notes}

Three readings that belong to particular statements above rather than to the comparison.

\begin{remark}[Pinning more than one row]\label{rem:multi}
Confining $j$ rows rather than one, each to $\sigma$ symbols, gives $R\bar c - 1 \ge j(W-1)/K$ by
the same averaging, provided a common support of size $t$ survives. It does: choosing the $j$ symbol
sets $\calS_1,\dots,\calS_j$ independently and uniformly among the subsets of $\calR$ of size
$\sigma$, each input survives all $j$ constraints with probability $(\sigma/R)^{j}$, so the expected
number of survivors is $D(\sigma/R)^{j}$ and some choice attains it; the constraint is
$D(\sigma/R)^{j} \ge t$, that is $j \le \log_{R/\sigma}(D/t)$. In the principal regime
$\sigma = 1$, $W = R$ that reads $j \le \log_R(D\epsilon)$, which is $14$, $6$ and $1$ in the three
rows of Section~\ref{sec:compare}. The gain is a factor $\sqrt{j}$ in the
bound, hence $\log_2 j$ further bits of seed in that comparison: at those three values of $j$,
between none and four. So the one-row construction of Theorem~\ref{thm:limit} carries most of the
effect, and none of it in the third row.

The averaging is the point. Choosing each $\calS_i$ greedily, as the $\sigma$ commonest symbols of
its row, does \emph{not} work for a worst-case table, since the $j$ greedy sets may have a small or
even empty common preimage; an earlier version of this material asserted the greedy version. Only the
one-row case survives greedily, and that is the case Theorem~\ref{thm:limit} uses.
\end{remark}

\begin{remark}[Where the loss comes from, and when it vanishes]\label{rem:cause}
The loss is the step \eqref{eq:cs-l1l2}, $\|v\|_1 \le \sqrt{KR}\,\|v\|_2$, which is tight only when
$v$ is spread evenly over all $KR$ cells; the source of Theorem~\ref{thm:limit} makes it far from
that, its deviation on the pinned row being a single concentrated lump.

The true advantage against that source should not be mis-stated, and an earlier version of this material
did mis-state it as $O(1/K)$. Since
$\SD(Q,U) = \frac1K\sum_k \SD(\mathrm{law}\,h(k,X), U_\calR)$, the pinned row $k_1$ contributes about
$1/K$, but the other $K-1$ rows are \emph{not} uniform. On a typical table each is $t$ draws into
$R$ bins and contributes $\Theta(\sqrt{R/t})$, so the true advantage is
$\Theta(\sqrt{R/t}) + \Theta(1/K)$, dominated by the \emph{unpinned} rows in the regime of
Remark~\ref{rem:numbers}. On an adversarial table, every row constant say, it is close to
$1$. So the overshoot of the lemma over the truth is neither $\sqrt{RK}$ nor uniform in the
parameters. In the instance of Remark~\ref{rem:numbers}, at the seed length \eqref{eq:conj}
requires, the lemma returns $\tfrac12\sqrt{(W-1)/K} = 2^{-6}$ against a true advantage of about
$0.4\sqrt{R/t} = 2^{-49.3}$; at the much larger seed length the lemma itself requires, the two come
within a small factor of one another. The floor is a statement about what the technique can certify,
not about how weak the source is.

The gap closes when $W$ does. At $\epsilon = 1/D$ one has $t = D$, $\sigma = R$, $W = 1$ and
Theorem~\ref{thm:limit} is trivial, correctly: a source of maximal entropy is uniform on all of
$\calD$ and has no support to select. More generally $W$ is small whenever the source's min-entropy
is close to $\log_2 D$ or the output is short: for $D = 2^{256}$, $R = 2^{64}$ and min-entropy $254$
one has $t = 2^{254}$, $\sigma = 2^{62}$, $W = 4$, and the theorem exhibits no loss at all.
\end{remark}

\begin{remark}[What this settles about the witness]\label{rem:qwitness}
Remark~\ref{rem:cause} decomposes the true advantage against that witness as
$\Theta(\sqrt{R/t}) + \Theta(1/K)$, the second summand from the pinned row and the first from the
$K-1$ unpinned ones, and notes that the first dominates. Theorem~\ref{thm:badcell} says where the
first comes from: it is an artefact of the distinguisher being unbounded. Against a $q_D$-bounded
distinguisher the unpinned rows contribute at most $q_D\epsilon = q_D/t$ in place of
$\sqrt{R/t}$, which is smaller exactly when $q_D \le \sqrt{Rt} = \sqrt{R/\epsilon}$, the threshold
of Corollary~\ref{cor:which} up to its factor of two, while the pinned row still contributes its
$s_S/K = 1/K$. So the reading that the advantage against this source is $\Theta(1/K)$ is correct precisely in the
query-bounded model and false outside it, which is what Remark~\ref{rem:cause} found the hard way.
\end{remark}

\section{Is the random oracle worse than universal hashing?}\label{sec:vs}

Proposition~\ref{prop:univ} gives a universal family the bound $\tfrac12\sqrt{\epsilon R}$ with no
second summand at all, while Theorem~\ref{thm:limit} saddles a random table with a floor of
$\tfrac12\sqrt{(W-1)/K}$ on top of it. Read quickly, that says the random oracle is the worse
extractor. The reading is wrong, and the error is a natural one.

\medskip
\noindent\textbf{The two are not being compared at the same seed length.} By
Proposition~\ref{prop:nouniv} a $2$-universal family on $D$ inputs needs $K \ge D/R$, a salt of
$n - m$ bits. That is what buys $\gamma \le 0$. The random-oracle question is posed at
$K \approx c^{2}\log_2 D/\delta^{2}$, a salt of $\log_2 n + 2\log_2(1/\delta)$ bits, which is
smaller by an exponential factor. Comparing the two bounds without comparing their seeds is
comparing the price of two objects while looking only at one price tag.

\medskip
\noindent\textbf{At equal seed length the random oracle matches, and the gap closes completely.}
Give the random table the salt a universal family requires, $K = D/R$. Then the conjectured second
summand is
\[
  \sqrt{\frac{\log_2 D}{K}} \;=\; \sqrt{\frac{R\log_2 D}{D}} ,
\]
which for $D = 2^{2048}$ and $R = 2^{128}$ is $2^{-954.5}$: not small, but absent. Both routes then
return $\Theta(\sqrt{\epsilon R})$, so there is nothing to choose between them. The second summand
is the price of a \emph{short} salt, not a defect of the random oracle. Universal hashing never
pays it only because it never has one.

\medskip
\noindent\textbf{At the seed length that makes the problem interesting, universal hashing does not
exist.} Take the first row of Section~\ref{sec:compare}'s table: a Diffie--Hellman element in a
$2048$-bit group, min-entropy $256$, a $128$-bit key, $\delta = 2^{-32}$. The conjectured bound
does the job with a $77$-bit salt. A universal family needs at least $1920$ bits, twenty-five times
as long, to do the same job no better: its error there is $\tfrac12\sqrt{\epsilon R} = 2^{-65}$
against the random oracle's $2^{-63}$ at the same salt. Measured by error per bit of seed, then,
universal hashing is the suboptimal construction and the random oracle the optimal one. This is
the standard picture. The leftover hash lemma is known to be far from optimal in seed length, and
\eqref{eq:conj} sits at the information-theoretic optimum, which universal hashing misses by an
exponential margin.

\medskip
\noindent\textbf{Why universality is immune to support selection is why it is expensive.}
Proposition~\ref{prop:univ} survives a source that reads the whole table, and its proof shows why:
the bound
$\bar\beta = \sum_{x \ne x'} p_x p_{x'}\Pr_k[h(k,x) = h(k,x')] \le (1-\CP)/R$ holds for
\emph{every} law $p$, because universality constrains \emph{every} pair $x \ne x'$ separately. A
worst-case-over-all-pairs hypothesis is automatically robust to an adversary who picks the pairs
last. But a worst-case-over-all-pairs hypothesis is exactly what Proposition~\ref{prop:nouniv}'s counting
argument charges $D/R$ seeds for. The immunity and the cost are the same property seen twice. A
random table satisfies only the average-case version, which is why choosing the support after
seeing $H$ can hurt it at all, and why $\log_2 n$ bits of salt suffice to stop that rather than
$n - m$.

\medskip
\noindent\textbf{And the floor is about the technique, not the object.} Theorem~\ref{thm:limit}
bounds what one accounting can certify about a random table. It is not a lower bound on the
advantage, and the random table is better than the floor suggests: \eqref{eq:conj} is
\emph{true}, proved in Part~I with $c = 2$, so the random oracle does achieve
$\sqrt{\epsilon R} + \sqrt{\log_2 D/K}$ at a short salt. Theorem~\ref{thm:limit} says only that the
collision route cannot prove it. Remark~\ref{rem:cause} says the same numerically: against the
witness source the true advantage is some $2^{-49}$ where the route returns $2^{-6}$.



\section{Prior art: the floor is Stinson's}\label{sec:prior}

Theorem~\ref{thm:limit} is not new. This section says so precisely.

Barak, Dodis, Krawczyk, Pereira, Pietrzak, Standaert and Yu~\cite{BDKPPSY11} open by naming two
limitations of the leftover hash lemma. The first is entropy loss. The second is the one at issue
here: for an (almost) universal family good enough for the LHL, the seed length must be at least
\[
  \min\bigl(u - v,\; v + 2\log(1/\varepsilon)\bigr) - O(1) ,
\]
where $u$ is the length of the source and $v$ the number of extracted bits, ``and, thus, must grow
with the number of extracted bits''. They attribute it to Stinson~\cite{Sti94}, and record
separately that Stinson showed perfectly universal families need a seed linear in $u$.

Translate: $u = \log_2 D$, $v = \log_2 R = m$, $\varepsilon = \delta$, and the seed length is
$\log_2 K$. That bound reads $\log_2 K \ge \min(\log_2 D - m,\ m + 2\log_2(1/\delta)) - O(1)$.
Corollary~\ref{cor:seed} gives $\log_2 K \ge \log_2(W-1) - 2 + 2\log_2(1/\delta)$, and in the
principal regime $W = R$ that is $m + 2\log_2(1/\delta) - 2$: the second branch. The first branch is
Proposition~\ref{prop:nouniv}'s $K \ge D/R$, which is Stinson's linear-in-$u$ result; the elementary
route given there reaches it only up to a factor of two, and the constant is Stinson's own. So the two salt columns of Section~\ref{sec:compare}'s table are the two branches of one bound
that has been known since 1994. And Appendix~\ref{sec:vs}'s reading, that universal hashing has the
longer salt rather than the shorter, is that bound's own reading. What Part~II adds is not the
number.

It adds three things, none of them a new bound. First, a derivation of the relevant branch from a
pigeonhole and two applications of Cauchy--Schwarz, with no probability in it, so the statement
holds for every table. Second, that derivation put into the game of Section~\ref{sec:setting}:
Theorem~\ref{thm:limit} gives a witness source per table, so the obstruction is an attack one can
write down rather than a counting bound on families. Third, Lemma~\ref{lem:attrib}, which separates
the cases where the floor speaks about $\gamma$ from those where the entropy deficiency explains it.
Whether any of the three is also in Stinson's papers I do not know: the attribution above
is~\cite{BDKPPSY11}'s, and I have not read~\cite{Sti94} at first hand.

Two things in the game of Section~\ref{sec:setting} that literature does \emph{not} cover.

\begin{itemize}
\item \textbf{The independence hypothesis is \emph{not} weakened here, and it is worth being exact
about that, because the natural summary is wrong.} Definition~2.1 of \cite{BDKPPSY11} quantifies
over all distributions $X$ with $\tilde H_\infty(X \mid Z) \ge m$, with the seed $S$ drawn
independently as the coins of a function $\mathrm{Ext}$ that is fixed beforehand and public. So a
source there may already depend on the \emph{family}; what it must not depend on is the
\emph{seed}. The game of Section~\ref{sec:setting} asks no more than that: $\sd$ is drawn after $S$ has
terminated, so its source is independent of the seed too. Nothing in the adversary model is
weakened, and Theorem~\ref{thm:limit}'s witness uses no freedom the leftover-hash setting withholds.
What differs is the \emph{family}: a uniformly random table in place of a universal one. The
content of Theorem~\ref{thm:limit} is that this substitution alone, with seed-independence fully
intact, already costs the seed length. (That paper's remark about not making ``the source $X$
dependent on the seed'' is about a different scenario, fixing one public seed and reusing it
indefinitely, and does not bear on this.)

\item \textbf{Their random-oracle comparison bounds one party only, and it is unsalted.} This is
set out in full in Remark~\ref{rem:roquery} below, because the asymmetry is the whole of why that
comparison does not reach this setting.
\end{itemize}

\begin{remark}[The query model of the random-oracle comparison]\label{rem:roquery}
Section~3.4 of \cite{BDKPPSY11} compares the leftover hash lemma against the ``dream'' bound
obtained by modelling a cryptographic hash as a random oracle. Its construction is \emph{unsalted}:
the derived key is $H(X)$, with no seed anywhere in the argument. The reasoning is that $H(X)$ is
distinguishable from uniform only if the attacker queries the oracle at $X$; as $X$ has min-entropy
$\mu$, that happens with probability at most $T/2^{\mu}$, where $T$ bounds the adversary's running
time. Now calibrate $T$ against the application's own security. A primitive $P$ that is
$\epsilon$-secure against complexity $T$ must satisfy $\epsilon \le T/2^{v}$, since exhaustive
search over a $v$-bit key must not succeed better than $\epsilon$, so $T \approx \epsilon 2^{v}$.
That gives
\[
  \epsilon_{\mathrm{RO}} \;\approx\; \frac{T}{2^{\mu}} \;=\; \epsilon\,2^{\,v-\mu} \;=\;
  \epsilon\,2^{-L} , \qquad L := \mu - v \ \text{the entropy loss} ,
\]
against $\sqrt{\varepsilon 2^{-L}}$ for their generalised LHL and $2^{-L/2}$ for the standard one
(the list in Section~\ref{sec:bounds} quotes all three).

\emph{Only one of the two parties is bounded, and the two are not bounded in the same sense.} The
\emph{distinguisher} is bounded, by the running time $T$, which in particular bounds its oracle
queries. The \emph{source} is not bounded, because it is not a query-making object at all. There
$X$ is a distribution of min-entropy $\mu$, fixed independently of $H$. No source algorithm with
oracle access appears in the model, so there is nothing about it to bound. Searching
\cite{BDKPPSY11} for query counts finds only the distinguisher's $T$ and the query budgets of the
\emph{applications} being keyed (encryption, MAC and PRF queries); the source never makes a query.

That asymmetry is exactly what the argument rests on, and it is not incidental. ``$H(X)$ is
distinguishable only if the attacker queried at $X$'' is false the moment $X$ may depend on $H$, and
Proposition~\ref{prop:unsalted} is how false: a source that reads $H$ and returns a uniform element
of its largest fibre is distinguished with advantage $1 - 1/R$ by an adversary making no query at
all, at a min-entropy of $\log_2 D - \log_2 R$. So this comparison needs the source independent of
$H$ \emph{itself}. That is strictly stronger than the leftover-hash half of the same paper needs,
where the source must be independent only of the seed and may know the family. Being unsalted, it has no seed to fall back
on, so independence from $H$ is the only thing holding it up. That is the hypothesis
the game of Section~\ref{sec:setting} does drop, and the salt is what it puts in its place.

There is also no formal random-oracle model to appeal to here: \cite{BDKPPSY11}'s Section~3.4 is
prose, reading ``it is instructive to compare\ldots the heuristic `dream' bound\ldots we claim'',
and the random oracle appears in no definition, lemma or theorem of that paper. Its formal
Definition~2.1 has $\mathrm{Ext}$ a plain function and $X$ a distribution quantified over, with no
oracle anywhere. So the question of how many oracle queries the source may make has no answer
there, not because the answer is zero but because the object is a distribution rather than an
algorithm.

It is also why that comparison needs no salt while the game of Section~\ref{sec:setting} does. Once
the distinguisher is unbounded the estimate $T/2^{\mu}$ is vacuous: in that game every party holds
all $KD$ entries of the table, so $T$ is effectively $KD$ and $T/2^{\mu} = KD\epsilon$ exceeds $1$
for every interesting $\epsilon$. Unsalted extraction against an unbounded distinguisher simply
fails, and this paper's own arithmetic agrees: at $K = 1$ the bound \eqref{eq:conj} reads
$c(\sqrt{\epsilon R} + \sqrt{\log_2 D})$, which is vacuous for every $D \ge 2$. The salt is what
rescues it, and the salt is what the question is about. Summarising the two models:
\begin{center}\small
\begin{tabular}{@{}lll@{}}
\toprule
 & \cite{BDKPPSY11}, Section~3.4 & this paper \\
\midrule
construction   & $H(X)$, unsalted                        & $H(\sd,x)$, salted \\
source         & a distribution, independent of $H$;     & an algorithm reading all of $H$; \\
               & makes no queries                        & unbounded \\
distinguisher  & running time $\le T$                    & unbounded; holds all of $H$ \\
bound          & $\epsilon\,2^{-L}$                      & $c(\sqrt{\epsilon R} + \sqrt{\log_2 D / K})$ \\
\bottomrule
\end{tabular}
\end{center}
\end{remark}


\begin{conjurabibliography}{9}

\bibitem{BDKPPSY11}
Boaz Barak, Yevgeniy Dodis, Hugo Krawczyk, Olivier Pereira, Krzysztof Pietrzak,
Fran\c{c}ois-Xavier Standaert and Yu Yu.
\newblock Leftover Hash Lemma, Revisited.
\newblock In \emph{Advances in Cryptology -- CRYPTO 2011}, LNCS 6841, pages 1--20. Springer, 2011.
\newblock Cryptology ePrint Archive, Report 2011/088, \url{https://eprint.iacr.org/2011/088}.
\newblock (Read at first hand, and quoted in Section~\ref{sec:bounds}: the seed-length and
entropy-loss limitations are the two bullets of its abstract, developed on page~1; the generalised
leftover hash lemma is its Theorem~3.1, applied in Theorems~3.3--3.7; expand-then-extract is its
Section~4; the unsalted random-oracle comparison is its Section~3.4, which is prose and carries no
theorem.)

\bibitem{McD89}
Colin McDiarmid.
\newblock On the method of bounded differences.
\newblock In \emph{Surveys in Combinatorics 1989}, London Mathematical Society Lecture Note
Series 141, pages 148--188. Cambridge University Press, 1989.
\newblock (Fact~\ref{fact:mcd} is the one-sided form of the \emph{independent bounded differences
inequality}, Lemma~(1.2) there, which is stated two-sidedly with a factor $2$; the one-sided bound
with the same constant is Theorem~(6.7), of which Lemma~(1.2) is the independent-coordinate
specialisation.)

\bibitem{RT00}
Jaikumar Radhakrishnan and Amnon Ta-Shma.
\newblock Bounds for dispersers, extractors, and depth-two superconcentrators.
\newblock \emph{SIAM Journal on Discrete Mathematics}, 13(1):2--24, 2000.
\newblock (The optimum $d \ge \log(n-k) + 2\log(1/\epsilon) - O(1)$ that \eqref{eq:conj} meets up
to $\log\frac{n}{n-k}$, since it pays $\log n$ where the optimum pays $\log(n-k)$; that is $0.2$
bits in the first row of Section~\ref{sec:compare}.
Read at first hand.)

\bibitem{Sti94}
Douglas R. Stinson.
\newblock Universal hashing and authentication codes.
\newblock \emph{Designs, Codes and Cryptography}, 4(4):369--380, 1994.
\newblock (Cited by~\cite{BDKPPSY11} as the source of that seed-length bound, and quoted in the
proof of Proposition~\ref{prop:nouniv} for the $\varepsilon$-almost-universal lower bound
$K \ge D(R-1)/\bigl(\varepsilon R(D-R) + R^{2} - D\bigr)$. \emph{Not} read at first hand for this
paper: the attribution is~\cite{BDKPPSY11}'s and the formula is taken from a secondary restatement.
Its special case $K \ge D/R$ at $\varepsilon = 1/R$ was checked independently here --- it is
attained by $h(k,(a,b)) = ak+b$ over $\mathbb{F}_q$, and a randomised search over $R \le 5$,
$K \le 9$ found no family beating it.)

\bibitem{Vad12}
Salil P. Vadhan.
\newblock \emph{Pseudorandomness}.
\newblock Foundations and Trends in Theoretical Computer Science. Chapter 6, Randomness Extractors.
\newblock (Theorem~6.17 is \eqref{eq:conj} in extractor parameters; the remark following
Theorem~6.18 is the seed length of universal hashing. Read at first hand.)

\end{conjurabibliography}

\end{document}
